\cleardoublepage
\phantomsection
\addcontentsline{toc}{part}{Calabi--Yau Quantum Groups}
\chapter*{Calabi--Yau Quantum Groups}
\begin{summarybox}
Extensions of geometric objects define Hall products. Tensor products
of representations require coproducts; exchange requires intertwiners.
A trace on a specified state space then gives an amplitude. We construct
these operations in categorical, surface and root models, and determine
the maps needed to join them to fields on a compact Calabi--Yau threefold.
\end{summarybox}

\chapter{Calabi--Yau structures and local observables}

\section{Calabi--Yau categories}
Let $\kfield$ be a field of characteristic zero and let $\Ccal$ be a
small stable $\kfield$-linear category. Smoothness means that its diagonal
bimodule is perfect; properness means that every morphism complex is
perfect over $\kfield$. These finiteness conditions lead to two cyclic
duality data. Write $CC_*(\Ccal)$ for Hochschild chains, with their circle
action, and $\Ccal^!$ for the inverse-dualizing bimodule.

A smooth, or left, $d$-Calabi--Yau structure is a negative-cyclic class
\[
 \theta_{\CY}:\kfield[d]\longrightarrow CC_*(\Ccal)^{hS^1}
\]
whose Hochschild image induces an equivalence
$\Ccal^![d]\simeq\Ccal$. A proper, or right, structure is a
circle-equivariant trace
\[
 CC_*(\Ccal)\longrightarrow\kfield[-d]
\]
whose composition pairing on morphisms is nondegenerate. For a smooth
proper category we use compatible left and right structures. A scalar
trace without its circle-equivariant lift is insufficient for this datum.

For an object $E$, the tangent complex of the derived moduli stack is
\[
 T_E\M_{\Ccal}\simeq\RHom_{\Ccal}(E,E)[1].
\]
Composition and the trace give its pairing.

\begin{theorem}[The two universal categorical outputs]
Let $\Ccal$ be smooth and proper with the compatible cyclic structures
above. Its Hochschild cochains carry a framed $\Etwo$-algebra structure.
If $\Ccal$ is of finite type and the moduli of objects satisfies the
representability hypotheses of the shifted-symplectic moduli theorem,
then $\M_{\Ccal}$ is $(2-d)$-shifted symplectic.
\end{theorem}
\begin{proof}
Cyclic Deligne applies to the nondegenerate circle-equivariant trace
and refines the Hochschild brace structure \cite{BravRozenblyum}.
The moduli pairing at $E$ is
\[
 \RHom(E,E)[1]^{\tensor2}\xrightarrow{\circ}\RHom(E,E)[2]
       \xrightarrow{\Tr}\kfield[2-d].
\]
Calabi--Yau duality makes it nondegenerate. The negative-cyclic lift
provides the closed form, globally on the representable moduli problem
\cite{PTVV,BravDyckerhoff}.
\end{proof}

The shift of this form depends on $d$. The Hochschild operations remain
framed $\Etwo$ operations. Neither output alone specifies a local field
theory or a factorization algebra on a chosen manifold.

\section{Local BV theories and factorization observables}
\subsection{Cyclic elliptic local Lie algebras}
Let $X$ be a complex manifold. A classical local field theory is specified by a sheaf $\Lcal$ of elliptic $L_\infty$ algebras, a local pairing inducing the degree-$-1$ BV form on $\Lcal[1]$, and an action functional $S$ satisfying the classical master equation. For an open set $U\subset X$, the classical observables are
\[
 \Obs^{\mathrm{cl}}(U)=C^*_{\CE,\mathrm{cont}}(\Lcal(U)).
\]
Here $\Lcal(U)[1]$ is the complex of unrestricted smooth fields. We use the
differentiable pro-cochain model and tensor products of the elliptic
observable construction \cite{CG1,CG2}. Continuous linear functionals
are compactly supported distributions. The underlying cochain product
is defined in this coefficient category; its Poisson bracket is defined
on the regular observable subcomplex on which the contractions exist.
An arbitrary pair of distributional functionals need not admit such a
contraction.
For $U\subset V$, restriction of fields $\Lcal(V)\to\Lcal(U)$
induces the covariant map on observables.
For pairwise disjoint $U_i\subset V$, restrict to
$\bigoplus_i\Lcal(U_i)$ and multiply the pulled-back cochains.
These are the prefactorization maps.

\begin{theorem}[Factorization observables]
An elliptic classical BV theory produces a factorization algebra $\Obs^{\mathrm{cl}}$. A renormalized quantization satisfying the quantum master equation produces a factorization algebra $\Obs^{\mathrm{q}}$ over $\kfield[[\hbar]]$.
\end{theorem}
\begin{proof}
Restriction is a morphism of local $L_\infty$ algebras, so its
contravariant continuous Chevalley map commutes with the differential.
For disjoint opens, the completed cochains of the finite direct sum
are the completed tensor product of the cochains. This constructs
the displayed maps and their composition identities.
The elliptic descent theorem in this completed coefficient category
supplies Weiss descent \cite{CG1,CG2}. In the quantum theory, renormalized time-ordered products define the structure maps, and the quantum master equation makes their differential square-zero and compatible with nested ultraviolet contractions \cite{CG1,CG2}.
\end{proof}

\subsection{Boundary and defect restriction}
Let a theory on $C\times\mathbb R$ be holomorphic on the curve and
locally constant on the oriented real line. A boundary or defect has its
own observable category; a normal half-line by itself does not equip
one fixed endomorphism complex with curve coefficients.

Choose instead a family of boundary objects $b_I$ on the curve's
configuration spaces, with diagonal transition maps and disjoint-locus
factorization equivalences. Suppose its derived endomorphism objects
and their compositions respect those maps. They define a coefficient
object $A_b$ with local values $\RHom(b_I,b_I)$.

\begin{theorem}[Ordered boundary algebra]
Under these coefficient and composition hypotheses,
\[
 A_b\in\Alg_{\Eone}(\Fact_D(C),\tensor^!).
\]
If each $b_I$ is a compact generator and the local module equivalences
preserve the specified coefficient transitions and factorization, they
descend to an equivalence with $A_b$-modules in these coefficients.
\end{theorem}
\begin{proof}
Composition and the identity give the internal associative operations.
Their compatibility with the coefficient maps is precisely the
condition that these operations are morphisms in $\Fact_D(C)$.
Compact-generator Morita theory gives the local module equivalences;
their assumed descent compatibility glues them.
\end{proof}

\section{Three curve-level algebraic outputs}
Let $C$ be a smooth curve and let $\Dcal_C$ be the fixed Ran category.  The same geometric source can produce three algebraic outputs, and their distinction is functorial.

\begin{definition}[Curve-level faces]
A transverse face is
\[
 A^\perp\in\Alg_{\Eone}(\Fact_D(C),\tensor^!).
\]
An associative chiral face is
\[
 A^{\ch}_{\mathsf{Ass}}\in\Alg_{\Eone}(\Dcal_C,\starCh).
\]
A chiral Lie face is
\[
 \g^{\ch}\in\Alg_{\mathsf{Lie}}(\Dcal_C,\starCh).
\]
Their bars are $\BarPerp$, $\BarAssCh$, and $\BarLieCh$.
\end{definition}

\begin{theorem}[Source separation]
The factorizing boundary construction gives the transverse face.
A multiplication for $\starCh$ gives the associative chiral face, and
its commutator gives the Lie face. In a specified stable symmetric
monoidal model, let $U_{\starCh}(\g)$ denote the associative enveloping
algebra of a Lie object. If the PBW and completed bar--Chevalley
comparison hold in that model, then
\[
 \BarAssCh(U_{\starCh}(\g))\simeq\BarLieCh(\g).
\]
A mixed construction requires a distributive law compatible with all
diagonal and factorization transitions.
\end{theorem}
\begin{proof}
The commutator is operadic restriction. The displayed comparison is
the enveloping bar--Chevalley map in the named monoidal category, with
the PBW filtration and convergence hypotheses just stated. A coherent
distributive law composes the two associative operations as in Volume~I.
\end{proof}

The associative envelope $U_{\starCh}$ is specified in this category;
it is not identified here with the Beilinson--Drinfeld chiral envelope.
Neither the dimension of the source nor an ordinary vector-space
calculation supplies the missing geometric coefficient maps.

\section{Holomorphic Chern--Simons theory}
Let $X$ be a complex threefold with holomorphic volume form $\Omega_X$, and let $\mathfrak a$ be a quadratic Lie algebra. The BV fields are
\[
 \mathcal E_X=\Omega^{0,*}(X,\mathfrak a)[1],
\]
with degree $-1$ pairing
\[
 \omega(A,B)=\int_X\Omega_X\wedge\angles{A,B}.
\]
Local integrals are evaluated on compactly supported variations. The action is
\[
 S_{\hCS}(A)=\int_X\Omega_X\wedge
 \left(\frac12\angles{A,\dbar A}+\frac16\angles{A,[A,A]}\right).
\]

\begin{theorem}[Classical hCS factorization algebra]
The action $S_{\hCS}$ satisfies the classical master equation. Its classical observables form a holomorphic factorization algebra on $X$.
\end{theorem}
\begin{proof}
The Hamiltonian vector field is the Maurer--Cartan differential $\dbar A+\frac12[A,A]$. Its square is the Dolbeault square, the derivation identity, and the Jacobi identity. The elliptic Dolbeault complex then gives factorization observables by the preceding theorem.
\end{proof}

A perturbative quantization is controlled by the local BV deformation complex. In complex dimension three, its one-loop gauge anomaly occupies the degree-four invariant-polynomial summand; mixed and gravitational components enter when the background fields carry the corresponding couplings. A quantization solves the quantum master equation order by order, trivializing its obstruction classes compatibly with restriction and factorization \cite{CostelloLiAnomaly}.

\section{Pushforward to a curve}
Let $p:X\to C$ be a holomorphic map to a complex curve. For a factorization algebra $\F$ on $X$, define
\[
 (p_*\F)(U)=\F(p^{-1}U).
\]

\begin{theorem}[Factorization pushforward]
The assignment $p_*\F$ is a factorization algebra on $C$. Suppose in addition that $p$ is proper and locally trivial in a topology compatible with the factorization theory, that $\F$ is locally constant in the fiber directions, and that product-chart excision and base change hold. Then on every trivializing open $U\subset C$ with fiber $F_p$ there is a natural equivalence
\[
 (p_*\F)(U)\simeq
 \int_{F_p}\F|_{F_p\times U},
\]
compatible on overlaps. To obtain algebraic factorization $D$-modules, require in addition
regular-holonomic algebraic models on every configuration space and
compatible algebraic models for all structure maps.
\end{theorem}
\begin{proof}
Disjoint opens in $C$ have disjoint inverse images, so the factorization products push forward. Weiss descent is inherited from inverse images of Weiss covers. On a product chart, local constancy in the fiber and factorization excision identify the value with fiberwise factorization homology. The base-change hypothesis glues these local equivalences. The additional algebraic models give the last realization.
Riemann--Hilbert can compare specified constructible and regular-holonomic
models; it does not algebraize an arbitrary holomorphic factorization algebra.
\end{proof}

Thus a curve shadow is attached to an actual map and an actual field theory. Different fibrations and boundary conditions produce different ordered chiral algebras from the same categorical source.

\section{The derived centre and the physical bulk}
For an ordered boundary algebra $A$, define
\[
 Z_{\ord}(A)=\RHom_{A\tensor A^{\op}}(A,A).
\]
The tensor and derived Hom are internal to the chosen closed monoidal
coefficient category, with admissible derived tensor products. The
Hochschild brace operations give its internal $\Etwo$ structure.
A coherent bulk action in a compatible coefficient model gives
\[
 \beta:\Obs_{\mathrm{bulk}}|_{\partial}\longrightarrow Z_{\ord}(A).
\]

\begin{definition}[Open--closed completeness]
A boundary condition is open--closed complete on a protected sector when $\beta$ is an equivalence on that sector.
\end{definition}

The definition turns the phrase ``bulk equals centre'' into a comparison theorem. The centre exists algebraically for every ordered boundary algebra; the physical bulk is identified with it by the open--closed map.

\chapter{Hall algebras and symplectic surfaces}

\section{Extension correspondences}
Let $\M=\coprod_{\gamma\in\Gamma}\M_\gamma$ be a moduli stack of objects in a fixed exact category or heart of a stable category. Let $\M^{\mathrm{ext}}_{\alpha,\beta}$ classify extensions
\[
 0\longrightarrow E_\beta\longrightarrow E\longrightarrow E_\alpha\longrightarrow0.
\]
The correspondence is
\[
 \M_\alpha\times\M_\beta
 \xleftarrow{\ p\ }
 \M^{\mathrm{ext}}_{\alpha,\beta}
 \xrightarrow{\ q\ }
 \M_{\alpha+\beta}.
\]

\begin{theorem}[Hall associativity]
Assume a monoidal coefficient theory with the required K\"unneth maps, pullbacks, pushforwards and coefficient comparisons, satisfying base change for the two-step filtration stack. Denote its resulting pull--push by $q_*p^*$. Then
\[
 a*b=q_*p^*(a\boxtimes b)
\]
defines an associative product.
\end{theorem}
\begin{proof}
The two parenthesizations are the two projections from the stack of filtrations $0\subset E_1\subset E_2\subset E_3$. Base change identifies both composites with pullback to this stack followed by pushforward to $E_3$.
\end{proof}

For critical CoHAs, the coefficient theory includes vanishing cycles and a coherent orientation. For symplectic CoHAs, it includes the selected equivariant Borel--Moore theory. Completion is chosen so that each output degree receives a convergent sum.

\section{Orientation data and the threefold Hall product}
An orientation is a square root of the determinant line of the
three-Calabi--Yau moduli problem. Strong orientation also specifies its
multiplicative isomorphisms on extension correspondences, compatible
with triple flags and the zero object. Joyce--Upmeier construct natural
orientations for the compact threefold moduli problems in their theorem;
the strong enhancement is a further condition \cite{JoyceUpmeier,KPS}.

Let $\Phi_\gamma$ be the Donaldson--Thomas perverse sheaf on
$\M_\gamma$. Define the shift $\kappa(\alpha,\beta)$ by the hyperbolic
comparison
\[
 (\Phi_\alpha\boxtimes\Phi_\beta)[\kappa(\alpha,\beta)]
       \simeq Rp_*q^!\Phi_{\alpha+\beta}.
\]
This fixes its sign convention. The index obeys
\[
 \kappa(\alpha,\beta)+\kappa(\alpha+\beta,\gamma)
 =\kappa(\beta,\gamma)+\kappa(\alpha,\beta+\gamma),
 \qquad \kappa(0,\alpha)=\kappa(\alpha,0)=0.
\]
Thus the tensor
\[
 (V\otimes_\kappa W)_\gamma
 =\bigoplus_{\alpha+\beta=\gamma}
       (V_\alpha\otimes W_\beta)[\kappa(\alpha,\beta)]
\]
has the induced associativity constraint.

\begin{theorem}[Threefold Hall algebra from strong orientation]
Assume the three-Calabi--Yau moduli stack satisfies the geometric
hypotheses of \cite[Theorem~8.6 and Corollary~8.8]{KPS}: its classical
stack is quasi-separated Artin with affine stabilizers, it is
$\Theta$-reductive, and the zero object is closed. Strong orientation
defines a unital associative product on
\[
 H=\bigoplus_\gamma H^\bullet(\M_\gamma,\Phi_\gamma),
 \qquad H\otimes_\kappa H\longrightarrow H.
\]
\end{theorem}
\begin{proof}
The hyperbolic comparison has adjoint
$p^*(\Phi_\alpha\boxtimes\Phi_\beta)[\kappa]\to
q^!\Phi_{\alpha+\beta}$. Properness of $q$ on the chosen filtration
component allows its counit to define a map on ordinary cohomology.
The triple-flag comparison and strong orientation identify the two
threefold products. The zero-endpoint correspondences give the unit
\cite[Proposition~8.29]{KPS}.
\end{proof}

\begin{definition}[Reduced $K3\times E$ Hall datum]
A reduced datum additionally specifies the elliptic translation
quotient, reduced critical coefficient complexes, their pull--push maps,
and orientation comparisons on all flags. A chosen support convention
and charge completion must make these maps defined and continuous.
A full quantum-group realization further requires the Hopf and pairing
data described below.
\end{definition}

The theorem uses ordinary critical cohomology. Neither an
$E$-equivariant Borel--Moore product nor descent of reduced coefficients
follows from strong orientation alone.

\section{The geometric \texorpdfstring{$E$}{E}-equivariant Hall source}
Let $S$ be a projective K3 surface, let $E$ be an elliptic curve, and put
\[
 X=S\times E.
\]
Write $\Coh_{\leq1}(X)$ for coherent sheaves whose support has dimension at most one.  This is an extension-closed abelian category.  Let $\Gamma_X$ be its numerical charge group and $\Gamma_X^+$ the
monoid of charges of coherent sheaves. The degree of the projected
one-cycle is additive:
\[
 m:\Gamma_X\longrightarrow\mathbb Z,\qquad
 m(\Gamma_X^+)\subseteq\mathbb Z_{\geq0}.
\]
For $\gamma\in\Gamma_X^+$, let $\M_\gamma$ be the moduli stack of objects of charge $\gamma$.  Translation on the elliptic factor gives an action of $E$ on every $\M_\gamma$ and on every extension stack.

Oberdieck and Shen construct the $E$-equivariant Grothendieck ring of stacks and its Hall convolution.  In the present notation the completed algebra is
\[
 \Hcal_X^{E,\mathrm{mot}}
 =\widehat\bigoplus_{\gamma\in\Gamma_X^+}
 K_0^E(\operatorname{St}/\M_\gamma),
\]
where only effective charges occur. Fix an admissible filtration for
which each coefficient of a product receives finitely many contributions.
It must control Euler characteristic and zero-dimensional charges as
well as curve class; a curve-degree filtration alone does not do this.

\begin{theorem}[Geometric reduced Hall algebra]
The extension correspondence defines the associative equivariant
motivic Hall product. Let $H_{\mathrm{reg}}$ be the Oberdieck--Shen
regular subalgebra and let
$H_{\mathrm{sc}}=H_{\mathrm{reg}}/(\mathbb L-1)H_{\mathrm{reg}}$ be
its semiclassical quotient, where $\mathbb L=[\mathbb A^1]$, with the coefficient localization of that construction.
Reduced integration gives a multiplicative map
\[
 I_X^E:\widehat H_{\mathrm{sc}}\longrightarrow
       \widehat{\mathbb Q[\Gamma_X]}[\epsilon]/(\epsilon^2)
\]
on every compatible admissible completion.
\end{theorem}
\begin{proof}
Translation preserves exact sequences and numerical charge. The
two-step flag correspondence therefore proves equivariant associativity.
The reduced integration theorem is \cite[Theorem~5]{OberdieckShen}.
For an equivariant variety $Z\to\M_\gamma$ its coefficient is
\[
 I_X^E[Z\to\M_\gamma]
 =\bigl(\chi(Z^E,\nu)-\chi((Z-Z^E)/E,\nu)\epsilon\bigr)e^\gamma,
\]
with the pulled-back Behrend function $\nu$ and weighted orbifold Euler
characteristic of the moving quotient. For sheaves supported in
dimension at most one on the threefold, Riemann--Roch gives
$\chi(F,G)=0$: both Chern characters begin in degree four, so their
product has no degree-six component. The target's Euler twist is thus
trivial. The theorem is multiplicative on $H_{\mathrm{sc}}$, and
continuity extends it to the selected completion.
\end{proof}

Its composite with $H_{\mathrm{reg}}\to H_{\mathrm{sc}}$ will also
be denoted $I_X^E$. It is not a map from every stack class with arbitrary
denominators. The sign of the moving term is part of the reduced convention.

\section{Local and wrapped geometry inside the Hall algebra}
Define the completed degree-zero and positive-degree pieces
\[
 \Hcal_X^{\mathrm{loc}}
 =\widehat\bigoplus_{m(\gamma)=0}
 K_0^E(\operatorname{St}/\M_\gamma),
 \qquad
 \Hcal_X^{\mathrm{wr}}
 =\widehat\bigoplus_{m(\gamma)>0}
 K_0^E(\operatorname{St}/\M_\gamma).
\]
A pure one-dimensional support of degree zero projects to a finite subset of $E$.  A support of positive degree projects onto $E$.  The first sector is point-local on the elliptic curve; the second is wrapped.

\begin{theorem}[Geometric local--wrapped decomposition]
Hall multiplication respects elliptic degree:
\[
 \Hcal_{m_1}*\Hcal_{m_2}\subseteq\Hcal_{m_1+m_2}.
\]
Consequently $\Hcal_X^{\mathrm{loc}}$ is a subalgebra, $\Hcal_X^{\mathrm{wr}}$ is a closed two-sided ideal, and
\[
 \Hcal_X^{E,\mathrm{mot}}
 \cong
 \Hcal_X^{\mathrm{loc}}\ltimes\Hcal_X^{\mathrm{wr}}
\]
as a complete vector space with multiplication determined by the Hall correspondences.  On its regular domain, reduced integration preserves this decomposition.
\end{theorem}
\begin{proof}
The cycle of the middle term of a short exact sequence is the sum of the cycles of its subobject and quotient.  Projection to $E$ therefore adds degrees.  Degree zero is closed under addition, and positive degree absorbs addition by a nonnegative degree.  The same grading is present on the target group algebra of $I_X^E$.
\end{proof}

This is the geometric local--wrapped grading. The root construction
has a separate grading of the same form. Comparing them requires a
primitive map that preserves this degree, the brackets and the coefficients.

\section{Cohomological refinement and the Gram character}
A cohomological refinement starts with a specified reduced equivariant
coefficient theory. Suppose it supplies complexes $\Phi^{\mathrm{red}}_\gamma$
on the chosen quotient or equivariant moduli stacks and admissible
pull--push maps on extension and flag correspondences. In a Borel--Moore
model write
\[
 \Hcal_X^{E,\mathrm{coh}}=
 \widehat\bigoplus_{\gamma\in\Gamma_X^+}
 H_*^{\mathrm{BM},E}(\M_\gamma,\Phi^{\mathrm{red}}_\gamma),
\]
including the degree shifts of that model.

\begin{theorem}[Cohomological lift from compatible coefficients]
If these maps satisfy the oriented base-change and triple-flag
identities, they define an associative product on
$\Hcal_X^{E,\mathrm{coh}}$. Elliptic degree gives a subalgebra in
degree zero and a closed ideal in positive degree.
\end{theorem}
\begin{proof}
The supplied coefficient identities make both threefold products the
same map through the flag correspondence. Additivity of numerical
charge gives the degree decomposition. The underlying correspondences
are the motivic ones; a realization between their coefficient theories
is further data.
\end{proof}

Use the Mukai lattice
$\widetilde H(S,\mathbb Z)=H^0(S,\mathbb Z)\oplus H^2(S,\mathbb Z)
\oplus H^4(S,\mathbb Z)$, with
\[
 ((r,c,s),(r',c',s'))_{\mathrm{Muk}}=c\cdot c'-rs'-r's.
\]
For the dyonic charge lattice
\[
 \Gamma_{\mathrm{phys}}
 =\widetilde H(S,\Z)\oplus\widetilde H(S,\Z),
\]
write
\[
 \Pi(Q,P)=\left(\frac{Q^2}{2},Q\cdot P,\frac{P^2}{2}\right).
\]
To compare sheaf charges with this lattice, choose an additive map
\[
 c:\Gamma_X\longrightarrow\Gamma_{\mathrm{phys}}.
\]
Put $B_\Pi(a,b)=\Pi(a+b)-\Pi(a)-\Pi(b)$ and give
$\widetilde\Gamma_{\mathrm{phys}}=\Gamma_{\mathrm{phys}}\times\mathbb Z^3$
the group law
\[
 (a,u)(b,v)=(a+b,u+v+B_\Pi(a,b)).
\]
Then $j_\Pi(a)=(a,\Pi(a))$ is a homomorphism.
Assume $j_\Pi c$ preserves the selected charge completions.
On the regular subalgebra $\mathcal H_{\mathrm{reg}}$ of the reduced
integration theorem, define
\[
 I_{X,\Pi,c}^E=(j_\Pi c)_*I_X^E:
 \mathcal H_{\mathrm{reg}}\longrightarrow
 \widehat{\mathbb Q[\widetilde\Gamma_{\mathrm{phys}}]}
 [\epsilon]/(\epsilon^2).
\]
Write $I_{X,\Pi}^E$ for this map when $c$ has been fixed.

\begin{proposition}[Polarized reduced character]\label{cyiv:polarized-character}
The map $I_{X,\Pi,c}^E$ is multiplicative.
Its coefficients are the reduced invariants supplied by $I_X^E$,
regraded by $j_\Pi c$.
If $g\in1+F^1H$ is group-like in a complete filtered Hall bialgebra
and the logarithm converges, then $\log g$ is primitive.
\end{proposition}
\begin{proof}
Expanding $\Pi(a+b+d)$ gives
\[
 B_\Pi(a,b)+B_\Pi(a+b,d)=B_\Pi(b,d)+B_\Pi(a,b+d).
\]
This proves associativity of the group law.
The definition gives $j_\Pi(a+b)=j_\Pi(a)j_\Pi(b)$.
Composing its group-algebra map with reduced integration proves
multiplicativity on $\mathcal H_{\mathrm{reg}}$.
For the last assertion, the two factors in
$\Delta g=(g\widehat\otimes1)(1\widehat\otimes g)$ commute.
Taking the convergent logarithm gives
$\Delta\log g=\log g\widehat\otimes1+1\widehat\otimes\log g$.
\end{proof}

The primitive reduced Gromov--Witten theorem supplies the separate
character $-\chi_{10}^{-1}$ in classes primitive in the K3 factor
\cite{OP}.
Comparing this character with the Hall character requires the same
charge map, sector, and normalization.
A primitive root framing must also preserve the extension bracket
and parity of the chosen BPS states.

\section{BPS primitives}
A coproduct defines primitives by
\[
 \Prim(H)=\{x:\Delta x=x\tensor1+1\tensor x\}.
\]
A BPS algebra can also be constructed geometrically before such a
coproduct is chosen. The distinction matters for the full CoHA.

\begin{theorem}[Two-Calabi--Yau BPS algebra]
In the setting of \cite{DHS}, assume the extension stack is proper over
the object stack, the moduli stack has the required three-term
resolution of derived Hom, its good moduli space is separated with
finite-type components, direct sum on that space is finite, and the
simple-object category has the specified right two-Calabi--Yau
duality. With the associative unital Hall product and the prescribed
sign twist, the BPS algebra is the positive enveloping algebra of the
constructed generalized Kac--Moody Lie algebra.
\end{theorem}
\begin{proof}
The simple BPS sheaves supply its generators. The extension product,
local quiver comparison and Serre relations give a morphism from their
generalized Kac--Moody envelope. The PBW and support decomposition in
\cite[Theorems~1.1--1.2]{DHS} identify this morphism with the BPS algebra.
\end{proof}

This theorem applies to the semistable K3 examples satisfying those
hypotheses. The full CoHA has a PBW description involving
$\mathfrak g_{\mathrm{BPS}}\otimes H^*(B\mathbb G_m)$; that description
is an isomorphism of the stated graded sheaves or vector spaces.
It does not identify the full Hall multiplication with multiplication
in $U(\mathfrak g_{\mathrm{BPS}})$.

\section{The \texorpdfstring{$A_2$}{A2} Hall algebra}
For the quiver $1\to2$ over $\mathbb F_q$, put $v^2=q$ and let
$E_i=[S_i]$ be the simple-object classes. Use the Hall twist
$v^{\langle\alpha,\beta\rangle}$, where
$\langle(a,b),(c,d)\rangle=ac+bd-ad$.
Counting subrepresentations gives
\[
 E_1^2E_2-(v+v^{-1})E_1E_2E_1+E_2E_1^2=0.
\]
For a representation of dimension vector $(2,1)$, the arrow has rank zero or one. The relevant flag counts are
\[
\begin{array}{c|ccc}
\rank&112&121&211\\ \hline
0&q+1&q+1&q+1\\
1&q+1&1&0.
\end{array}
\]
After the Hall twists $v^{-1},1,v$, both rows give zero. This is the elementary geometric origin of the quantum Serre relation \cite{Ringel,Green}.

\section{The tripled Jordan quiver and \texorpdfstring{$\mathbb C^3$}{C3}}
For the one-vertex quiver with loops $x,y,z$, take the potential
$W=\operatorname{Tr}x[y,z]$. Its Jacobi algebra is $\mathbb C[x,y,z]$.
Dimensional reduction compares its critical CoHA with the CoHA of
zero-dimensional sheaves on $\mathbb A^2$.

\begin{theorem}[Affine-Yangian positive half]
Let $T=(\mathbb C^*)^2$ act with weights $t_1,t_2$ on the plane and
weight $-t_1-t_2$ on the third loop. Over
$H_T^*(\mathrm{pt},\mathbb Q)=\mathbb Q[t_1,t_2]$, the surface CoHA
is isomorphic to the strictly positive affine Yangian of
$\mathfrak{gl}_1$ in the presentation of \cite{DavisonAffineBPS}.
The isomorphism sends its current generator $e_i$ to the corresponding
one-point class $\alpha_1^{(i)}$.
\end{theorem}
\begin{proof}
Dimensional reduction identifies the Hall products. Davison proves
spherical generation and the current relations over the polynomial
parameter ring, then uses the affine BPS description and PBW to prove
the isomorphism. The shuffle comparison gives an additional localized
presentation with $\epsilon_1+\epsilon_2+\epsilon_3=0$
\cite{DavisonAffineBPS,DavisonMeinhardt,SVYangians,Tsymbaliuk}.
\end{proof}

No parameter localization is needed for the stated polynomial-ring
isomorphism. Nonequivariantly the algebra is
$U(\operatorname{gr}_F W_{1+\infty}^+)$; the one-parameter version is
the corresponding Rees envelope. Here
\[
 W_{1+\infty}^+=\operatorname{span}_{\mathbb Q}
   \{z^mD^n:m>0,\ n\geq0\},\qquad D=z\frac{d}{dz},
\]
and
\[
 [z^mD^a,z^nD^b]
 =z^{m+n}\bigl((D+n)^aD^b-D^a(D+m)^b\bigr).
\]
The order filtration is
$F_i=\operatorname{span}\{z^mD^a:a\leq(i+2)/2\}$.
The full affine Yangian has separate negative and Cartan currents.
A vacuum representation yields a vertex algebra only after a vacuum,
translation operator, fields and a reconstruction theorem have been
specified; an unspecified evaluation homomorphism does not provide them.

\section{Surface CoHAs and ADE Yangians}
Let $S_\Gamma\to\mathbb C^2/\Gamma$ be the minimal resolution for
a finite subgroup $\Gamma\subset\operatorname{SL}_2(\mathbb C)$,
and let $C$ be its exceptional curve. Use the diagonal torus action
commuting with $\Gamma$.

\begin{theorem}[Curve-supported surface CoHA]
The equivariant CoHA of coherent sheaves set-theoretically supported
on $C$, including their zero-dimensional subobjects, has the completed
nonstandard affine ADE Yangian presentation of \cite{DPPSSV}.
Its completion uses quasi-compact and slope bounds; its Yangian
transition maps are the truncated braid maps between quotients of
the standard negative half.
\end{theorem}
\begin{proof}
The reflection equivalences and Hecke correspondences identify the
finite slope pieces. Their compatibility with the truncated braid maps
gives the inverse-system comparison. The multiplication and continuity
theorems in \cite{DPPSSV} give the completed isomorphism.
\end{proof}

An ADE configuration on a K3 surface suggests a local application.
It requires an identification of the supported categories, coefficient
theories and equivariant geometry. The presence of the intersection
matrix alone does not supply that identification or a global torus action.

\section{Coulomb branches and shifted Yangians}
For a three-dimensional $\mathcal N=4$ gauge theory, the BFN Coulomb branch is defined by equivariant Borel--Moore homology of an affine-Grassmannian correspondence. Its convolution algebra quantizes the Coulomb branch \cite{BFN}. In quiver gauge theories, shifted Yangians and their truncations provide presentations of these quantizations. In models with a specified boundary construction, its geometric data determine the relevant shift and truncation.

The BFN construction supplies a second geometric source of Yangians. Loop rotation supplies the quantization parameter and convolution
supplies the product. A coproduct requires an additional construction.
The surface CoHA uses sheaf extensions. A comparison between the two is a model-specific equivalence.

\section{The global zero-dimensional surface algebra}
Let $S$ be a projective K3 surface, and let $\M_S(n)$ classify
zero-dimensional sheaves of length $n$. In the cohomological grading
and rational Borel--Moore theory of \cite{MMSVSurface}, put
\[
 H_S^0=\bigoplus_{n\geq0}H_*^{\mathrm{BM}}(\M_S(n),\mathbb Q).
\]
Choose a hyperbolic plane with basis $e,f$ in $H^2(S,\mathbb Z)$,
where $e^2=f^2=0$ and $e\cdot f=1$. The class
$q=e-12f$ satisfies
$q^2=-24[\mathrm{pt}]=c_1(S)^2-c_2(S)$.
This is a cohomology class; no Hodge-type condition on $q$ is required.

Define $\mathfrak w_S^+$ on the symbols
$z^mD^n\lambda$, linear in $\lambda\in H^*(S,\mathbb Q)$, with
$m>0$, $n\geq0$, by
\[
 [z^mD^n\lambda,z^aD^b\mu]
 =z^{m+a}\frac{(D+aq)^nD^b-D^n(D+mq)^b}{q}\,\lambda\mu.
\]
The quotient means the polynomial obtained before specializing the
formal parameter to $q$. Its numerator is divisible by that parameter;
the nilpotent class $q$ is never inverted. To see Jacobi, first impose
$Dz=z(D+t)$ over $\mathbb Q[t]$. Its commutator divided by $t$ is a
polynomial Lie bracket, so specialization $t\mapsto q$ preserves Jacobi.

Write $W^+(S)$ for the positive surface current algebra. It is
generated by $T_n(\lambda)$, $n\geq0$, linearly in
$\lambda\in H^*(S,\mathbb Q)$, with coefficient migration
\[
 [T_m(\lambda\mu),T_n(\nu)]
   =[T_m(\lambda),T_n(\mu\nu)],
\]
the current relation
\[
\begin{aligned}
 0={}&[T_m(\lambda),T_{n+3}(\mu)]
 -3[T_{m+1}(\lambda),T_{n+2}(\mu)]\\
 &+3[T_{m+2}(\lambda),T_{n+1}(\mu)]
 -[T_{m+3}(\lambda),T_n(\mu)]\\
 &-[T_m(\lambda),T_{n+1}(q^2\mu)]
 +[T_{m+1}(\lambda),T_n(q^2\mu)],
\end{aligned}
\]
and the cubic relation
\[
 \sum_{\sigma\in S_3}\sigma\cdot
 [T_{m_3}(\lambda_3),[T_{m_2}(\lambda_2),T_{m_1+1}(\lambda_1)]]=0.
\]
Here the indices are nonnegative and $\sigma$ permutes the three
pairs $(m_i,\lambda_i)$.
\begin{theorem}[Global surface CoHA]
There is an algebra isomorphism
\[
 H_S^0\simeq W^+(S)\simeq U(\mathfrak w_S^+).
\]
The length-one current class indexed by $D^n\lambda$ corresponds to
$zD^n\lambda$.
\end{theorem}
\begin{proof}
The surface CoHA theorem identifies $H_S^0$ with the positive
deformed $W$-algebra \cite[Theorem~B]{MMSVSurface}. For $c_1=0$ and
$c_1^2-c_2=q^2$, the current relations identify this algebra with
$U(\mathfrak w_S^+)$ \cite[Theorem~3.5]{MMSVSurface}. The proof
constructs the current homomorphism; surjectivity and the proved PBW
dimensions make it an isomorphism.
\end{proof}

\section{A global surface coproduct}
Let $\Gamma$ be the lattice of Mukai vectors of coherent sheaves on
$S$, where $v(F)=\operatorname{ch}(F)\sqrt{\operatorname{td}(S)}$.
Let $\chi$ be the surface Euler pairing, the negative of the Mukai
pairing above. For the stack $\M_S(\gamma)$ of objects of vector
$\gamma$, put $v_\gamma=-\chi(\gamma,\gamma)$, its virtual dimension.
Work in rational
pro-complexes obtained by restriction to quasi-compact open substacks:
\[
 A_\gamma=R\Gamma^{\mathrm{BM}}_{\mathrm{pro}}
          (\M_S(\gamma),\mathbb Q)[-v_\gamma].
\]
Write $A=(A_\gamma)_\gamma$ for this charge-graded pro-object. There are two charge tensors. On total charge $\gamma$ they are
\[
\begin{aligned}
 (A\widehat\otimes'_{\Gamma}B)_\gamma
 &=\bigoplus_{\alpha+\beta=\gamma}A_\alpha\widehat\otimes B_\beta,\\
 (A\widehat\otimes_{\Gamma}B)_\gamma
 &=\prod_{\alpha+\beta=\gamma}A_\alpha\widehat\otimes B_\beta.
\end{aligned}
\]
Here each tensor on the right is the tensor of the specified pro-objects.
In addition to the cochain Koszul sign, the symmetry has the factor
$(-1)^{\chi(\alpha,\beta)}$.

The coefficient construction uses oriented shifted Lagrangian classes
\[
 [L]_{\mathrm{Lag}}:f_1^*\Phi_1\longrightarrow
             f_2^!\Phi_2[-\operatorname{vdim}L]
\]
for a correspondence $X_1\xleftarrow{f_1}L\xrightarrow{f_2}X_2$.
We require their unit, composition, external-product and \'etale-base-change
identities, together with dimensional reduction for shifted cotangent
stacks. These are the functorial coefficient hypotheses used in
\cite[Theorem~2.4 and Section~6]{DHKSvSurface}.

\begin{theorem}[Eigenvalue coproduct for a symplectic surface]
With this coefficient package, the surface Hall construction gives
\[
 m:A\widehat\otimes'_{\Gamma}A\longrightarrow A,
 \qquad
 \Delta:A\longrightarrow A\widehat\otimes_{\Gamma}A.
\]
The coproduct is coassociative and cocommutative for the Euler-twisted
symmetry. It is compatible with the Hall product in the completed
tensor domains of \cite[Proposition~6.4]{DHKSvSurface}. The resulting
cohomology is a topological graded bialgebra.
\end{theorem}
\begin{proof}
Enhance a sheaf to $(F,\phi)$ with $\phi\in\operatorname{End}(F)$.
For an open $U\subset\mathbb C$, restrict to endomorphisms whose
eigenvalues lie in $U$. Dimensional reduction identifies the critical
cohomology over a nonempty convex $U$ with $A$. For two disjoint such
opens, generalized eigenspaces split the sheaf canonically. Oriented
Thom--Sebastiani identifies the corresponding critical complexes with
the product-completed tensor of their two factors.

Restriction from a convex open containing both gives $\Delta$.
Three disjoint opens give the same restriction in either order,
proving coassociativity. Interchanging the two summands changes the
orientation by $(-1)^\chi$, which gives the stated cocommutativity.
The zero-sheaf component gives the counit.

For multiplicativity, split each term of an extension into its two
generalized-eigenvalue parts. Off-diagonal extensions between disjoint
spectra vanish. The resulting square of extension correspondences
identifies the two product--coproduct composites. The composition and
external-product identities for the shifted Lagrangian classes make
their coefficient maps agree; the middle interchange has the Euler
and Koszul signs. Properness of the flag maps gives the locally finite
extension of multiplication on this completed target
\cite[Propositions~6.1 and~6.4]{DHKSvSurface}.
\end{proof}

For this argument the object stack is quasi-smooth, locally of finite
type with affine diagonal, closed under direct sums and summands,
contains the zero object as an open and closed component, and is
$\Theta$-reductive. The pro-system handles its non-quasi-compactness.
Multiplication is not asserted on an arbitrary product over all charge
decompositions.

Let $H_S=\bigoplus_\gamma H^*(A_\gamma)$ after topological realization.
Fix a polarization and a reduced Hilbert polynomial $p$, and write
$H_S^{\mathrm{ss}}(p)$ for the corresponding semistable Hall algebra. The semistable
charge cone is strictly convex and its cohomology has finite graded
pieces. Let $\mathfrak g_S$ and $\mathfrak g_S^{\mathrm{ss}}(p)$ be
the primitives for the two coproducts. In this same symmetry,
\[
 U(\mathfrak g_S^{\mathrm{ss}}(p))\simeq H_S^{\mathrm{ss}}(p),
 \qquad U(\mathfrak g_S)\hookrightarrow H_S
\]
with dense image in the second map. Restriction identifies the primitives
of each semistable charge with the corresponding global primitives;
their graded spaces are the affinized BPS cohomology, including the
$H^*(B\mathbb G_m)$ factor and the normalization of dimensional reduction
\cite[Theorem~6.8]{DHKSvSurface}.

The semistable assertion follows from connected cocommutative
reconstruction. For the global assertion, ordered Harder--Narasimhan
products give PBW tensors of semistable primitive envelopes. An iterated
coproduct followed by semistable restriction detects the first nonzero
Harder--Narasimhan term, proving injectivity. The same triangular
argument approximates every class on each quasi-compact truncation,
proving density. Purity and the strict-nilpotent-support comparison
identify the primitive spaces in this argument.

This constructs a global K3 coproduct on the stated surface coefficients.
Its comparison with reduced critical coefficients on $K3\times E$ is
a separate problem. Its pro-tensors and Euler symmetry also differ
from the fixed-parameter polynomial shuffle carrier whose ordinary
Koszul cooperation fails in Volume~I.

\chapter{Amalgams, doubles and tensor representations}

\section{The integral Hall--Borcherds amalgam}
Let $\mathfrak p_X^+$ be a complete charge-graded BPS Lie superalgebra
and let $\nDelta$ be the positive Igusa Borcherds Lie superalgebra.
Choose a finite-dimensional abelian Cartan algebra $\mathfrak h_\Pi$ and additive
weight maps from both charge groups to $\mathfrak h_\Pi^*$.
For a homogeneous $x_\gamma$, the action is
$[h,x_\gamma]=\lambda(\gamma)(h)x_\gamma$.
Additivity of $\lambda$ makes this action a derivation. Form
\[
 \mathfrak b_X=\mathfrak h_\Pi\ltimes\mathfrak p_X^+,
 \qquad
 \mathfrak b_\Delta=\mathfrak h_\Pi\ltimes\nDelta.
\]
The two chosen weights must agree along any primitive comparison.
The polarized group law makes $j_\Pi$ additive, but its coordinate
$\Pi(a)$ is not an additive scalar weight.

For example, let $[x,y]=z$ in the Heisenberg Lie algebra, with charges
$1,1,2$. An additive grading derivation has eigenvalues $1,1,2$.
The quadratic prescription $D(x)=x$, $D(y)=y$, $D(z)=4z$ gives
\[
 D[x,y]=4z,\qquad [Dx,y]+[x,Dy]=2z.
\]
It therefore cannot define a Cartan action. The polarized character
and the additive Cartan action remain separate constructions.

Fix a nonnegative filtration whose positive parts have strictly positive
integer height and finite-dimensional pieces. Keep the Cartan in degree
zero, with a specified compatible coefficient topology. Complete the
ordinary amalgam by this filtration; all mapping properties below are
in the resulting category of continuous filtered maps.

\begin{definition}[Integral protected Lie algebra]
The integral Hall--Borcherds algebra is the completed amalgamated coproduct
\[
 \mathfrak b_{X\boxtimes\Delta}
 =\mathfrak b_X
  \widehat\amalg_{\mathfrak h_\Pi}
  \mathfrak b_\Delta
\]
in complete charge-filtered dg Lie superalgebras.
\end{definition}

\begin{theorem}[Hall--Borcherds amalgamation]
The completed amalgam exists and represents pairs of continuous Lie morphisms from $\mathfrak b_X$ and $\mathfrak b_\Delta$ that agree on $\mathfrak h_\Pi$.  The Cartan projections of the two Borels induce canonical retractions
\[
 r_X:\mathfrak b_{X\boxtimes\Delta}\longrightarrow\mathfrak b_X,
 \qquad
 r_\Delta:\mathfrak b_{X\boxtimes\Delta}\longrightarrow\mathfrak b_\Delta.
\]
Its completed enveloping algebra is the amalgamated free product
\[
 \widehat U(\mathfrak b_{X\boxtimes\Delta})
 \cong
 \widehat U(\mathfrak b_X)
 \widehat *_{\widehat U(\mathfrak h_\Pi)}
 \widehat U(\mathfrak b_\Delta).
\]
The two factor inclusions are split injections. Their Lie brackets and
enveloping products are retained. A representation or character of
either factor pulls back along its retraction; this does not identify
the two factors' characters or construct a compact determinant map.
\end{theorem}
\begin{proof}
Construct the ordinary Lie coproduct from the free product of the two Borels and divide by the ideal identifying their two Cartan copies.  Complete by the chosen positive-height filtration. Finite-dimensional
graded pieces and positivity make each fixed-height Lie calculation
finite. Degree-zero Cartan brackets reduce to the prescribed weight
action. The separated inverse limit has the stated continuous mapping
property.  The projection $\mathfrak b_X\to\mathfrak h_\Pi\to\mathfrak b_\Delta$ together with the identity on $\mathfrak b_\Delta$ gives $r_\Delta$ by universality; the second retraction is symmetric.  The enveloping functor preserves coproducts of Lie algebras, and PBW identifies the completed algebraic filtration with the displayed completed free product.
\end{proof}

\begin{definition}[Comparison ideal]
A primitive comparison $\varphi:\nDelta\to\mathfrak p_X^+$ preserving Cartan degree defines the closed ideal
\[
 \mathcal J_\varphi
 =\overline{\angles{
   i_\Delta(x)-i_X(\varphi(x))\mid x\in\nDelta
 }}
 \subset\mathfrak b_{X\boxtimes\Delta}.
\]
The diagonal realization is
\[
 \mathfrak b_{X\simeq\Delta,\varphi}
 =\mathfrak b_{X\boxtimes\Delta}/\mathcal J_\varphi.
\]
\end{definition}

\begin{proposition}[Graph quotient of a primitive comparison]\label{cyiv:graph-quotient}
Let $\varphi:\mathfrak b_\Delta\to\mathfrak b_X$ be a continuous Lie
morphism restricting to the identity on $\mathfrak h_\Pi$.
Use a completion in which the amalgam and closed quotient have
their stated universal properties.
Then
\[
 \mathfrak b_{X\boxtimes\Delta}/\mathcal J_\varphi
 \cong\mathfrak b_X.
\]
The two factor maps become $\id$ and $\varphi$.
If $\varphi$ is an isomorphism, the quotient is also isomorphic
to $\mathfrak b_\Delta$.
\end{proposition}
\begin{proof}
The maps $\id$ and $\varphi$ agree on $\mathfrak h_\Pi$.
They give $q_\varphi:\mathfrak b_{X\boxtimes\Delta}\to\mathfrak b_X$,
which kills the generators of $\mathcal J_\varphi$.
The inclusion of $\mathfrak b_X$ gives a map $s$ in the other direction.
The composite $q_\varphi s$ is the identity.
The other composite is the identity on the geometric factor.
On the Borcherds factor, the defining relation identifies $x$ with
$s\varphi(x)$, so it is also the identity there.
The two factor images topologically generate the quotient.
Continuity proves that the maps are inverse on the completion.
\end{proof}

The graph quotient implements a specified primitive comparison.
Constructing that comparison requires geometric simple states whose
extension brackets satisfy the Borcherds relations.
The quotient formula alone produces neither those states nor their
orientation data.

\section{Coproducts, pairings, and doubles}
A finite double shows which structures enter before completion.
Let $G$ be a finite group and let $D(G)$ have basis
$\delta_g h$, $g,h\in G$. Write $e$ for the group identity and set
\[
 (\delta_g h)(\delta_x k)
 =\begin{cases}\delta_g(hk),&g=hxh^{-1},\\0,&\text{otherwise},\end{cases}
 \qquad 1=\sum_g\delta_g e.
\]
Its coalgebra and antipode are
\[
 \Delta(\delta_g h)=\sum_{ab=g}\delta_a h\otimes\delta_b h,
 \quad \epsilon(\delta_g h)=\delta_{g,e},
 \quad S(\delta_g h)=\delta_{h^{-1}g^{-1}h}h^{-1}.
\]
Put $1h=\sum_x\delta_xh$.

\begin{theorem}[Finite group double]
These formulas make $D(G)$ a Hopf algebra with universal matrix
\[
 R=\sum_g\delta_g e\otimes1g,
 \qquad R^{-1}=\sum_g\delta_g e\otimes1g^{-1}.
\]
It satisfies
\[
 R\Delta(a)=\Delta^{\mathrm{op}}(a)R,\quad
 (\Delta\otimes\id)R=R_{13}R_{23},\quad
 (\id\otimes\Delta)R=R_{13}R_{12}.
\]
\end{theorem}
\begin{proof}
A triple product is nonzero precisely when
$g=hxh^{-1}$ and $x=kyk^{-1}$; either order then gives
$\delta_g(hkl)$. This proves associativity. Both iterated coproducts
sum over $abc=g$. Conjugation respects multiplication in $G$, so
$\Delta$ preserves products. The two counit identities follow by
setting one factor of $ab=g$ equal to $e$.

In $m(S\otimes\id)\Delta(\delta_gh)$ a term survives only when
$b=a^{-1}$, hence only when $g=e$. In that case the sum is
$\sum_a\delta_{h^{-1}a^{-1}h}e=1$. The other antipode identity is
the same calculation. Orthogonality of the $\delta_g$ gives the
displayed inverse of $R$.

Write $R$ in its function and group factors. Moving a group element
past a function conjugates its index, which gives
$R\Delta(\delta_gh)=\Delta^{\mathrm{op}}(\delta_gh)R$.
For the first hexagon, both sides sum the terms with indices $a,b,ab$.
For the second, orthogonality in the first tensor factor forces the
same index $g$ in both remaining group factors. These are the stated
two equations. Applying either counit to $R$ gives $1$.
\end{proof}

The conjugation and the inverse in $S$ are essential. Already for
$G=S_3$, erasing conjugation leaves the tensor-product algebra
$\mathbb Q^G\otimes\mathbb QG$, but the displayed $R$ no longer
intertwines its coproduct. Erasing $g^{-1}$ in $S$ destroys the
convolution identity.

For a finite-dimensional Hopf algebra $H$ with invertible antipode,
the convention of \cite[Section~7.14]{EGNOTensor} pairs $H$ with $H^{*\mathrm{cop}}$.
The opposite coproduct and antipode enter the cross-relations; the
dual bases give the canonical tensor. This gives a full positive model
without any convergence assumption.

\begin{definition}[Completed Hall double]
Choose positive and negative Hopf Borel halves with invertible
antipodes, their common Cartan identifications, and a Hopf pairing with
the specified opposite-coproduct convention. Require a separated
triangular multiplication in the chosen completed tensor category,
with underlying module
\[
 \widehat D(H^+)=H^-\widehattensor H^0\widehattensor H^+.
\]
Here the notation includes the chosen negative and Cartan data. Its
canonical tensor must be supplied by locally perfect dual pieces or
by an explicit continuous coevaluation, and must converge on the
selected modules. The Cartan reduction must be a Hopf quotient.
This datum defines the completed Hall double.
\end{definition}

\begin{theorem}[Double construction]
Suppose the finite dual pieces or finite quotients in this datum carry
their compatible ordinary doubles, cross-relations and canonical
tensors. Their separated limit is a complete Hopf double. On the
module domain where its canonical tensor converges, it is
quasitriangular and hence defines braided tensor products.
\end{theorem}
\begin{proof}
At each finite stage the Hopf pairing and the opposite dual coproduct
give the ordinary Drinfeld double. Compatibility makes every
associativity, antipode and quasitriangular equation an equation of
inverse systems. Passing to the separated limit preserves those
equations in the specified tensor domains. The coevaluation hypothesis
supplies the limiting canonical tensor; its inverse and the two
hexagons are inherited from the finite stages. A compatible Cartan
Hopf quotient preserves these identities
\cite{DrinfeldDouble,EtingofSchiffmann,EGNOTensor}.
\end{proof}

Nondegeneracy of a pairing alone need not produce a coevaluation or
a convergent sum of dual bases. A general infinite Hall construction
must establish those conditions directly when no finite-double system
is available.

\begin{theorem}[Quantum Hall--Borcherds amalgam]
Let $H_X$ and $H_\Delta$ be complete Hopf algebras containing the same complete Cartan Hopf algebra $H^0$, with continuous Hopf retractions onto $H^0$.  Form their algebraic amalgamated free product and choose a separated
multiplicative filtration for which the two coproducts and antipodes
extend continuously to the completed tensor products. Their completed
amalgamated free product
\[
 H_{X\boxtimes\Delta}
 =H_X\widehat *_{H^0}H_\Delta
\]
has a unique complete Hopf structure extending the two coproducts, counits, and antipodes.  It represents the pairs of filtered continuous Hopf morphisms whose induced map is continuous for this filtration, and has Hopf retractions to the factors when those retractions preserve the filtration.  A universal $R$-matrix requires an invertible element in its completed
tensor square satisfying intertwining with both generating factors,
the two hexagons and the counit normalizations. The two factor
$R$-matrices do not by themselves supply it.
\end{theorem}
\begin{proof}
The ordinary amalgam is the free product modulo the relations
identifying the Cartan copies; its chosen separated completion exists
as their filtered inverse limit.  The two coproducts agree on $H^0$, so the universal property extends them to a coproduct on the free product.  Coassociativity is checked on the two generating Hopf subalgebras.  The same argument extends the counit and antipode.  The Hopf retractions are induced by the filtration-compatible factor retractions onto $H^0$.  All extensions are continuous by the filtration hypothesis.
The additional quasitriangular equations must hold on both generating
factors and in the chosen tensor completion.
\end{proof}

A retraction from a full quantum group to its unrestricted Cartan is
itself a substantial assumption. For $U_q(\mathfrak{sl}_2)$ over
$\mathbb Q(q)$, killing $E,F$ in
$[E,F]=(K-K^{-1})/(q-q^{-1})$ forces $K^2=1$. The usual positive
Borel projection therefore does not extend to that Cartan retraction
of the full quantum group.

The deconcatenation coproduct on a bar complex and the representation-theoretic coproduct on a quantum group occupy different structures. A geometric comparison between them is a theorem about a selected model.

\section{Ordered chiral bialgebras}
A tensor product of representations requires compatible actions on
threefold products and the unit.
Retain the factorization coefficient category
$(\Fact_D(C),\tensor^!)$ on the smooth complex curve $C$.
Fix a category $\widehat{\mathcal V}_C$ of complete filtered objects
with specified presentations $F\simeq\varprojlim_r F_r$,
where $F_r\in\Fact_D(C)$.
Use derived limits in the fixed Ran enhancement and define
\[
 F\widehat{\tensor}^{!}G
 =\varprojlim_{r,s}(F_r\tensor^!G_s).
\]
Assume these limits remain factorizing and give
$\widehat{\mathcal V}_C$ a symmetric monoidal structure with unit $\mathbf1$.
Its associativity, symmetry, and unit equivalences must preserve
the presentations and the Ran structure.
Morphisms are continuous maps compatible with these presentations.
These are hypotheses on the chosen completion.

\begin{definition}[Ordered chiral bialgebra]\label{def:vol4-ordered-chiral-bialgebra}
An ordered chiral bialgebra is a counital $\Eone$-coalgebra object
in complete ordered chiral algebras:
\[
 A\in\operatorname{CoAlg}_{\Eone}
 \left(\Alg_{\Eone}
 (\widehat{\mathcal V}_C,\widehat{\tensor}^{!})\right).
\]
It has continuous multiplication $m$ and unit $\eta$, and continuous
morphisms of unital ordered chiral algebras
\[
 \Delta:A\longrightarrow A\widehat{\tensor}^{!}A,
 \qquad \epsilon:A\longrightarrow\mathbf1.
\]
The counit is the chosen augmentation.
Write $\tau$ for the coefficient symmetry, including Koszul signs.
Suppressing ambient associativity and unit equivalences, compatibility reads
\[
\begin{aligned}
 \Delta m&=(m\widehat{\tensor}^{!}m)
  (\id\widehat{\tensor}^{!}\tau\widehat{\tensor}^{!}\id)
  (\Delta\widehat{\tensor}^{!}\Delta),\\
 \Delta\eta&=\eta\widehat{\tensor}^{!}\eta,
 \qquad \epsilon m=\epsilon\widehat{\tensor}^{!}\epsilon,
 \qquad \epsilon\eta=\id_{\mathbf1}.
\end{aligned}
\]
Coassociativity and the counit identities read
\[
\begin{aligned}
 (\Delta\widehat{\tensor}^{!}\id)\Delta
   &=(\id\widehat{\tensor}^{!}\Delta)\Delta,\\
 (\epsilon\widehat{\tensor}^{!}\id)\Delta
   &=\id_A=(\id\widehat{\tensor}^{!}\epsilon)\Delta.
\end{aligned}
\]
For a homotopy-coherent presentation, these identities include the
higher coherences of the displayed coalgebra object.
Every operation and coherence preserves diagonal transitions and
disjoint-locus factorization maps.
An antipode is a continuous convolution inverse $S:A\to A$ of $\id_A$:
\[
 m(S\widehat{\tensor}^{!}\id)\Delta
 =\eta\epsilon
 =m(\id\widehat{\tensor}^{!}S)\Delta.
\]
A bialgebra with an antipode is a Hopf object in this coefficient category.
\end{definition}

A continuous left $A$-module is an object $M\in\widehat{\mathcal V}_C$
with a unital $\Eone$ action
$\rho_M:A\widehat{\tensor}^{!}M\to M$.
Its action and coherences are morphisms in these factorization coefficients.
Write $\operatorname{Mod}_A(\widehat{\mathcal V}_C)$ for their category.

\begin{theorem}[Representation fusion]\label{thm:vol4-representation-fusion}
An ordered chiral bialgebra makes
$\operatorname{Mod}_A(\widehat{\mathcal V}_C)$ monoidal.
Its product is $M\widehat{\tensor}^{!}N$ with action
\[
\begin{aligned}
 A\widehat{\tensor}^{!}(M\widehat{\tensor}^{!}N)
 &\xrightarrow{\Delta\widehat{\tensor}^{!}\id}
 (A\widehat{\tensor}^{!}A)\widehat{\tensor}^{!}
 (M\widehat{\tensor}^{!}N)\\
 &\xrightarrow{\id\widehat{\tensor}^{!}\tau_{A,M}
                       \widehat{\tensor}^{!}\id}
 (A\widehat{\tensor}^{!}M)\widehat{\tensor}^{!}
 (A\widehat{\tensor}^{!}N)\\
 &\xrightarrow{\rho_M\widehat{\tensor}^{!}\rho_N}
 M\widehat{\tensor}^{!}N.
\end{aligned}
\]
The unit is $\mathbf1$ with action induced by $\epsilon$.
The ambient associator and unit equivalences are module equivalences.

Suppose a continuous universal $R$-matrix is also given.
It is an invertible morphism
$R:\mathbf1\to A\widehat{\tensor}^{!}A$ for the tensor-product algebra,
satisfying
\[
\begin{gathered}
 R\Delta(a)=\Delta^{\op}(a)R,\qquad
 \Delta^{\op}=\tau_{A,A}\Delta,\\
 (\Delta\widehat{\tensor}^{!}\id)(R)=R_{13}R_{23},
 \qquad
 (\id\widehat{\tensor}^{!}\Delta)(R)=R_{13}R_{12},\\
 (\epsilon\widehat{\tensor}^{!}\id)(R)=\eta
 =(\id\widehat{\tensor}^{!}\epsilon)(R).
\end{gathered}
\]
The first equation is an identity of morphisms from $A$ to its completed tensor square.
Leg notation uses the unit and symmetry of $\widehat{\mathcal V}_C$.
In a homotopy-coherent presentation, require compatible lifts of these identities.
Then $c_{M,N}=\tau_{M,N}\circ R_{M,N}$ is a braiding,
where $R_{M,N}$ acts through the external tensor-product module action.
\end{theorem}
\begin{proof}
The symmetry combines the independent module actions into an
$A\widehat{\tensor}^{!}A$-action.
Restriction along the unital algebra map $\Delta$ gives the displayed action.
For three modules, the two actions use
$(\Delta\widehat{\tensor}^{!}\id)\Delta$ and
$(\id\widehat{\tensor}^{!}\Delta)\Delta$.
Coassociativity makes the ambient associator $A$-linear.
The counit identities make the ambient unit equivalences $A$-linear.
The coalgebra coherences give the pentagon and triangle coherences.
The completion hypotheses keep these maps continuous and factorizing.
Quasitriangularity makes $c_{M,N}$ an intertwiner and invertibility makes it an equivalence.
The two coproduct identities for $R$ give the hexagons.
Its counit identities give compatibility with the unit.
\end{proof}

The tensor in this construction is the completed transverse tensor
$\tensor^!$ on factorization coefficients.
A chiral-convolution or Hall coproduct requires its own construction
and a comparison in the specified category.
A spectral version additionally requires a difference coordinate,
a Laurent-series completion, and compatible domains for iterated expansions.
Its identities must hold in those domains before defining module actions.

\begin{example}[A counital algebra map without coassociativity]
In ordinary complex vector spaces, let $A=\C[x]$ and put
\[
 \Delta(x)=x\tensor1+1\tensor x+x\tensor x^2,
 \qquad \epsilon(x)=0.
\]
Polynomial substitution extends these assignments to unital algebra maps.
Writing $F(u,v)=u+v+uv^2$, the equations
$F(0,v)=v$ and $F(u,0)=u$ prove both counit identities.
But
\[
\begin{split}
 F(F(u,v),w)-F(u,F(v,w))
 &=-uv^2w^4-uv^2w^2-2uvw^3-2uvw\\
 &=-uvw(1+w^2)(vw+2)\ne0.
\end{split}
\]
On three regular modules, the two actions differ already on
$1\tensor1\tensor1$.
Thus the canonical tensor associator is not $A$-linear.
This ordinary example shows why multiplicativity and the counit identities
must be supplemented by coassociativity.
It makes no assertion about a geometric realization in factorization coefficients.
\end{example}

\section{Tensor reconstruction and geometric kernels}
Let $\mathscr C$ be an essentially small rigid linear monoidal category
and let $\omega:\mathscr C\to\operatorname{Vect}_{\mathrm{fd}}(\mathbb Q)$
be a $\mathbb Q$-linear strong monoidal functor.
A choice of geometric kernels specifies the morphisms of $\mathscr C$;
it is part of the input.
Choose a small skeleton and define
\[
 H=\left(\bigoplus_{M\in\mathscr C}\omega(M)^*\otimes\omega(M)\right)
 \big/\left\langle
 [\lambda\circ\omega(f)\otimes v]_M
 -[\lambda\otimes\omega(f)v]_N\right\rangle,
\]
where $f:M\to N$, $v\in\omega(M)$, and $\lambda\in\omega(N)^*$.

\begin{proposition}[Matrix-coefficient coalgebra]\label{cyiv:k3-coend}
For a basis $(e_i)$ with dual basis $(e^i)$, the formulas
\[
 \Delta[\lambda\otimes v]_M
 =\sum_i[\lambda\otimes e_i]_M\otimes[e^i\otimes v]_M,
 \qquad \epsilon[\lambda\otimes v]_M=\lambda(v)
\]
are basis-independent and make $H$ a coalgebra.
Tensor product gives its multiplication, and duals give its antipode.
Each $\omega(M)$ is a comodule by
\[
 v\longmapsto\sum_i e_i\otimes[e^i\otimes v]_M.
\]
Every selected kernel is a comodule morphism.
\end{proposition}
\begin{proof}
The identity tensor $\sum_i e_i\otimes e^i$ is basis-independent.
Inserting it twice proves coassociativity; evaluation proves both
counit identities.
The defining relations express naturality under $f$ and are preserved
by these formulas.
The strong monoidal constraints identify the coefficient of $M\otimes N$
with the product of coefficients of $M$ and $N$.
Their pentagon and unit identities give the bialgebra axioms.
The evaluation and coevaluation identities for duals give the antipode
identities.
The same relations prove that $\omega(f)$ intertwines the coactions.
\end{proof}

The linear dual of $H$ represents linear natural endomorphisms of
$\omega$.
Tensor-compatible natural automorphisms are the corresponding
multiplicative, unit-preserving families, not all elements of this dual.
An equivalence with a comodule category requires the additional exactness
and reconstruction hypotheses stated in the Tannaka theorem below.
An action algebra generated by selected kernels retains those
operators; this coend encodes the natural coefficient coactions that
make the selected morphisms intertwiners.

\section{Tannaka reconstruction}
Let $\mathcal L$ be an essentially small rigid $\kfield$-linear
abelian monoidal category with finite-dimensional morphisms,
finite-length objects and exact tensor product. Let
\[
 \omega:\mathcal L\longrightarrow\Vect_{\mathrm{fd}}
\]
be faithful, exact and strong monoidal. Form its matrix-coefficient
Hopf algebra $H$ by the preceding coend construction. The underlying
vector space of $H$ may be infinite-dimensional.

Extend $\omega$ to
$\overline\omega:\operatorname{Ind}(\mathcal L)\to\Vect$.
Assume it has a right adjoint $R$, is comonadic, and that the canonical
coend map identifies its comonad with $-\otimes H$. Require also that
objects with finite-dimensional $\overline\omega$ are exactly the
objects of $\mathcal L$.

\begin{theorem}[Rigid Tannaka reconstruction]
Under these hypotheses the comparison is a monoidal equivalence
\[
 \mathcal L\simeq\operatorname{Comod}_{\mathrm{fd}}(H).
\]
A braiding on $\mathcal L$ gives the corresponding coquasitriangular
form on $H$.
\end{theorem}
\begin{proof}
The coend formulas give each $\omega(M)$ its universal coaction.
Tensor products and duals give the Hopf operations. Comonadicity
identifies $\operatorname{Ind}(\mathcal L)$ with coalgebras for
$\overline\omega R\simeq-\otimes H$. Restricting to finite-dimensional
underlying spaces gives the displayed equivalence. The two hexagons
transport to the two multiplicativity equations of the
coquasitriangular form.
\end{proof}

The comonad lives on $\Vect$; an infinite-dimensional $H$ does not
define an endofunctor $-\otimes H$ of $\Vect_{\mathrm{fd}}$.
A spectral version additionally needs a parameter base, meromorphic
coefficient domains and compatible expansions. Ordinary rigid
reconstruction does not supply those data.

\section{From positive half to full quantum group}
A local or compact positive algebra enters a full quantum group
through Hopf Borel halves, a continuous pairing, Cartan reduction and
the tensor completion of the double construction. On the modules
where the canonical tensor acts, it gives braided tensor products.
For a topologically free Hopf deformation $H_\hbar$ of $U(\mathfrak g)$,
the first antisymmetric term
\[
 \delta(x)=\hbar^{-1}(\Delta-\Delta^{\mathrm{op}})(\widetilde x)
                         \pmod\hbar
\]
defines its semiclassical cobracket on primitives, when the stated
lift and topological freeness conditions hold.

The affine Yangian of $\mathfrak{gl}_1$ has independently presented
negative and Cartan currents. The exceptional-curve algebra has its
specified nonstandard completed positive presentation. The root
Borcherds construction in Volume~III has a paired Borel realization
with its own topology. None of these statements identifies an
arbitrary reduced compact Hall source with a full quantum group.

\section{KZ and rational exchange}
For a simple complex Lie algebra $\mathfrak g$, fix an invariant
nondegenerate form and let $\Omega$ be its inverse tensor. On a tensor
product of modules, the Knizhnik--Zamolodchikov connection is
\[
 \nabla=d-\frac{\hbar}{2\pi i}
 \sum_{i<j}\Omega_{ij}\,d\log(z_i-z_j).
\]
Flatness follows from the infinitesimal braid relations. Over formal
$\mathbb C[[\hbar]]$, the Drinfeld--Kohno comparison identifies the
corresponding categories of deformations of finite-dimensional modules
after a choice of associator, with $q=e^{\hbar/2}$ in this normalization
\cite[Section~8.27.7]{EGNOTensor}. A specialization requires its own
parameter and convergence hypotheses.

For the standard module of $\mathfrak{gl}_N$, rational difference exchange uses the Yangian matrix
\[
 R(u)=1-\frac{\hbar P}{u}
\]
where $P(v\otimes w)=w\otimes v$. The Yang--Baxter equation is the confluence identity for three exchanges. The KZ and Yangian constructions share the geometry of configuration spaces while using different flat connections.

\section{The centre and the double}
The Hochschild centre $Z_{\ord}(A)$ controls central operations on the
ordered algebra in its specified coefficient category.
To form a Drinfeld centre of modules, first specify a monoidal module category.
For an ordered chiral bialgebra satisfying
Definition~\ref{def:vol4-ordered-chiral-bialgebra}, use
\[
 \mathcal M_A=\operatorname{Mod}_A(\widehat{\mathcal V}_C)
\]
with the monoidal product of Theorem~\ref{thm:vol4-representation-fusion}.
An object of $Z(\mathcal M_A)$ is a module $M$ together with natural equivalences
\[
 c_{M,N}:M\widehat{\tensor}^{!}N
 \xrightarrow{\sim}N\widehat{\tensor}^{!}M
 \qquad(N\in\mathcal M_A)
\]
satisfying the half-braiding hexagon and unit identities, with their coherences.
This categorical centre is braided.
A comparison with the Hochschild centre must specify a realization functor,
the source and target of the comparison, and the operations it preserves.
Dualizability alone does not supply such a comparison.
Derived loop-space descriptions apply in the geometric categories and
under the hypotheses of their respective centre theorems
\cite{BenZviFrancisNadler}.

The centre and the double answer complementary questions.
The centre classifies central actions in the chosen category.
The double reconstructs a quantum group from paired positive and negative halves.
A field-theory comparison requires explicit realizations of the relevant
line and bulk operations and a map compatible with both constructions.

\chapter{Geometric actions on K3 moduli}

\section{The Mukai lattice}
For a K3 surface $S$, the Mukai lattice is
\[
 \Mukai(S,\Z)=H^0(S,\Z)\oplus H^2(S,\Z)\oplus H^4(S,\Z)
\]
with pairing
\[
 ((r,c,s),(r',c',s'))=c\cdot c'-rs'-r's.
\]
It is even unimodular of signature $(4,20)$ and rank $24$.
The Heisenberg commutator below uses the Poincar\'e pairing
$\int_S\alpha\beta$, which has different signs on $H^0\oplus H^4$.
These two pairings are not identified without that sign change.

\section{The Nakajima Heisenberg algebra}
Let
\[
 \mathbb H_S=\bigoplus_{n\ge0}H^*(S^{[n]}).
\]
Use rational cohomology. For $m\in\Z\setminus\{0\}$ and $\alpha\in H^*(S)$, normalize the Nakajima operators $p_m(\alpha)$ so positive modes annihilate the vacuum and the following commutator holds.

\begin{theorem}[K3 Heisenberg action]
The operators satisfy
\[
 [p_m(\alpha),p_n(\beta)]
 =m\,\delta_{m+n,0}\int_S\alpha\cup\beta\;\id.
\]
They generate the rank-twenty-four Heisenberg algebra acting irreducibly on $\mathbb H_S$ with vacuum $1\in H^*(S^{[0]})$.
\end{theorem}
\begin{proof}
The commutator is the excess intersection of the two incidence correspondences. Opposite creation and annihilation correspondences meet along the diagonal with multiplicity $m$ and intersection pairing $\int_S\alpha\beta$. Correspondences of the same sign commute. The vacuum construction gives the Fock representation \cite{Nakajima97}.
\end{proof}

This is a nonabelian central extension: opposite modes have the
displayed scalar commutator. Each creation half is abelian. Its Fock
representation is a global K3 construction; a larger action must
contain it through an actual homomorphism.

\section{The orthogonal LLV algebra}
For a compact hyperkähler manifold, every Lefschetz class gives an $\mathfrak{sl}_2$ triple on cohomology. The Lie algebra generated by all such triples is the Looijenga--Lunts--Verbitsky algebra.

\begin{theorem}[Orthogonal LLV symmetries]
Over $\mathbb R$, for a K3 surface $S$, the LLV algebra acting on
$H^*(S,\mathbb R)$ is
\[
 \mathfrak g_{\mathrm{LLV}}(S)
 \cong\mathfrak{so}\bigl(\widetilde H(S),(-,-)_{\mathrm{Muk}}\bigr)
 \cong\mathfrak{so}(4,20).
\]
For an irreducible hyperk\"ahler manifold $M$ of $K3^{[n]}$ type with $n\ge2$, one has $b_2(M)=23$ and
\[
 \mathfrak g_{\mathrm{LLV}}(M)\cong\mathfrak{so}(4,21).
\]
Smooth compact hyperk\"ahler moduli spaces of stable objects on $S$ carry their own LLV actions; Mukai and monodromy correspondences relate these actions to the K3 lattice.
\end{theorem}
\begin{proof}
The Looijenga--Lunts--Verbitsky theorem identifies the Lie algebra generated by all Lefschetz triples with the orthogonal algebra of the Mukai extension $H^2(M)\oplus U$. For $S$, the extension has signature $(4,20)$. For $K3^{[n]}$ type with $n\ge2$, the Beauville--Bogomolov lattice has signature $(3,20)$, so adjoining the hyperbolic plane gives signature $(4,21)$ \cite{LooijengaLunts,Verbitsky96}.
\end{proof}

Over $\mathbb Q$ the orthogonal Lie algebra is specified by its
rational quadratic form, rather than by a real signature alone.
The LLV algebra is a classical cohomological symmetry. A Yangian realization of the full orthogonal action requires a current algebra, its action on the stated cohomology, and a compatible spectral coproduct.

\section{Local ADE Yangians}
An ADE configuration of $(-2)$-curves on $S$ gives a root sublattice of $H^2(S,\Z)$. An identification of the supported categories and equivariant
coefficient systems with the resolution model gives its completed
positive affine ADE Yangian. Such an identification is additional
geometry, beyond the intersection matrix. Nakajima quiver varieties supply compatible Yangian actions on their equivariant cohomology.

Thus the ADE models provide local nonabelian Yangian sectors when
these geometric comparisons are given. Picard monodromy transports the ADE root sublattices over the corresponding Noether--Lefschetz loci. A family of sheaf moduli and Hecke correspondences compatible with that monodromy transports the Yangian actions themselves.

\section{Global and local K3 algebras}
The global zero-dimensional CoHA is $W^+(S)$, and it admits the
associative double $W^{(c)}(S)$ of \cite{MMSVSurface}. This algebra has
positive and negative current halves and a commutative Cartan part,
with a linear triangular decomposition
\[
 W^-(S)\otimes W^0(S)\otimes W^+(S)\simeq W^{(c)}(S).
\]
At an integer central parameter $c=r$ the algebra is defined over
$\mathbb Q$. Its positive currents $T_n^+(\lambda)$ have bidegree
$(1,2n-2+\deg\lambda)$; negative currents have first degree $-1$.
The algebra is defined by the current presentation, including its mixed relation, in \cite[Definition~3.10]{MMSVSurface}.

The Heisenberg embedding has
\[
 [q_m(\lambda),q_n(\mu)]
 =m\delta_{m+n,0}\,c\!\int_S\lambda\mu.
\]
For $c_1(S)=0$, the colored Virasoro central form is
$\eta(\lambda,\mu)=c\int_Sc_2(S)\lambda\mu$; in the convention
of \cite[Proposition~3.20]{MMSVSurface},
\[
 [L_m(\lambda),L_n(\mu)]
 =(n-m)L_{m+n}(\lambda\mu)
  -\frac{n^3-n}{12}\delta_{m+n,0}\eta(\lambda,\mu).
\]
This is the opposite of the usual Virasoro sign convention. On a K3,
$\eta(1,1)=24c$.

A two-sided Hecke pattern is a family of moduli spaces, open in the
specified sheaf stacks, stable under the elementary modifications used
for both creation and annihilation. With the regularity and properness
hypotheses of \cite[Proposition~6.8]{MMSVSurface}, a rank-$r$ pattern
carries an action of $W^{(r)}(S)$ on its tautological classes.
The virtual-class formulation gives the corresponding action without
the regularity assumption in that theorem. The Hilbert-scheme case
recovers the Nakajima operators. Separately, the universal
Fock-type module of \cite[Proposition~4.8]{MMSVSurface} is faithful
for $r>0$. It is neither irreducible nor a highest-weight module;
the geometric Hilbert-scheme representation can factor through a
central quotient of the double.

These are global associative algebra and representation results.
They do not assert a Hopf double or a universal $R$-matrix for
$W^{(c)}(S)$, nor an embedding of the full LLV algebra into it.
The following comparisons retain their separate carriers:
\begin{center}
\small
\begin{tabularx}{0.98\textwidth}{@{}p{0.25\textwidth}XX@{}}
\toprule
Structure&Carrier&Operation\\\midrule
Heisenberg&$\bigoplus_nH^*(S^{[n]})$&creation, annihilation\\
LLV&$H^*(S)$ or $H^*(M)$ for hyperk\"ahler $M$&Lefschetz commutators\\
Global $W^{(r)}(S)$&specified two-sided Hecke patterns&positive and negative currents\\
Local ADE Yangian&equivariant exceptional-curve moduli&completed Hecke currents\\
Affinized BPS algebra&surface Hall pro-cohomology&Euler-twisted coproduct\\
\bottomrule
\end{tabularx}
\end{center}
Maps between these actions must compare their state spaces, coefficients
and mixed commutators.

\section{K3 correspondences on fixed coefficients}
Fix rational cohomology and a collection $\mathscr M_S$ of smooth
projective moduli spaces on a projective K3 surface $S$.
Include the Hilbert schemes $S^{[n]}$ and close the collection under
finite products.
For $M,N\in\mathscr M_S$, define the graded kernel space
\[
 K(M,N)=H^*(M\times N,\mathbb Q)[2\dim_{\mathbb C}M].
\]
A kernel of degree $r$ acts from $H^j(M)$ to $H^{j+r}(N)$ by
\[
 T_\alpha(v)=(p_N)_*(p_M^*v\cup\alpha).
\]
The shift in $K(M,N)$ accounts for integration over $M$.
All pushforwards in this construction are proper.

\begin{proposition}[Kernel composition]\label{cyiv:k3-kernels}
The product
\[
 \beta\circ\alpha=(p_{13})_*
 (p_{12}^*\alpha\cup p_{23}^*\beta)
\]
defines an associative graded category with diagonal identities.
The kernel action identifies $K(M,N)$ with
$\operatorname{Hom}^*(H^*(M),H^*(N))$.
\end{proposition}
\begin{proof}
The K\"unneth isomorphism and Poincar\'e duality identify
$H^*(M\times N)[2\dim M]$ with $H^*(M)^\vee\otimes H^*(N)$.
Under this identification, $T_\alpha$ is the associated linear map.
Integration over the middle factor contracts its dual pair, so
$T_{\beta\circ\alpha}=T_\beta T_\alpha$.
Faithfulness then proves associativity and the diagonal identity.
External products give tensor products with their Koszul signs.
\end{proof}

A geometric symmetry is specified by selected kernels in this category.
Nakajima incidence kernels connect different Hilbert schemes.
Lefschetz multiplication and its adjoint act on each individual
cohomology space.
Using every cohomology kernel would instead give all linear maps;
the selected geometric kernels carry the representation-theoretic content.

\section{The algebra generated by geometric actions}
Put $V_S=\bigoplus_{M\in\mathscr M_S}H^*(M,\mathbb Q)$.
Choose a set of homogeneous operators $T_\lambda$ on $V_S$, each sending
a vector with finite support to a vector with finite support.
The Nakajima operators on the Hilbert-scheme summand satisfy this
condition: each changes the point number by a fixed integer.
Individual Lefschetz operators, extended by zero on other summands,
also satisfy it.
Let $F$ be the free unital graded algebra on symbols $t_\lambda$.

\begin{proposition}[Geometric action algebra]\label{cyiv:k3-action}
The map
\[
 \rho:F\longrightarrow\operatorname{End}_{\mathbb Q}(V_S),
 \qquad t_\lambda\longmapsto T_\lambda
\]
defines a faithful action of $A_S=F/\ker\rho$.
For every algebra $B$, maps $A_S\to B$ are exactly maps $F\to B$
that kill $\ker\rho$.
Products and commutators of the selected kernels are their convolution
products and commutators on the relevant moduli spaces.
\end{proposition}
\begin{proof}
Finite words in the selected operators preserve finite support, so
$\rho$ is defined by composition.
The quotient acts faithfully by definition of its kernel.
The factorization property is the universal property of an algebra
quotient.
The convolution assertion follows from Proposition~\ref{cyiv:k3-kernels}.
\end{proof}

An equivariant ADE construction uses a separate coefficient field,
for example $\operatorname{Frac}H_T^*(\mathrm{pt},\mathbb Q)$,
and its prescribed localization and proper pushforwards.
Joining it to $A_S$ requires maps between the actual state spaces that
intertwine those operators and their coefficients.
A completion of $A_S$ requires a specified topology and continuity of
every generator.
The indefinite Mukai norm alone supplies neither condition.
The desired common K3 quantum algebra must construct these local-to-global
actions and their mixed convolution relations.

\section{The three K3 brackets}
The Heisenberg bracket is the excess intersection of Nakajima
incidence kernels. The LLV bracket is the commutator of Lefschetz
kernels and their adjoints. Write $\operatorname{Corr}_{K3}(S)$ for
the rational kernel category constructed above. An ADE Yangian
initially lives in its equivariant coefficient category. When an actual
coefficient and state comparison puts its operators on a common
carrier, mixed commutators are differences of the corresponding
triple pushforwards, with Koszul signs. The action algebra retains
these operators and relations. The coend additionally requires the
specified rigid monoidal category and fibre functor.

\chapter{Protected roots, cyclic operations and characters}

\section{Finite quivers and Calabi--Yau completions}
For a finite quiver $Q$ with zero potential, let $\Gamma(Q,0)$ be its
path-length-completed Ginzburg algebra.
Path products are read from left to right: if $a:i\to j$ and
$b:j\to k$, then $ab:i\to k$.
For each arrow $a:i\to j$, add $a^*:j\to i$ of degree $-1$.
At each vertex add a loop $t_i$ of degree $-2$.
Original arrows have degree zero, and
\[
 da=da^*=0,\qquad
 dt_i=e_i\left(\sum_a[a,a^*]\right)e_i.
\]

\begin{proposition}[Restriction to a full subquiver]\label{cyiv:ginzburg-projection}
If $Q_0$ is a full subquiver, retaining its generators and killing all
other vertices, arrows, reverse arrows, and loops defines a continuous
dg-algebra projection
\[
 \Gamma(Q,0)\longrightarrow\Gamma(Q_0,0).
\]
The inclusion of retained graded generators need not commute with
the differential.
\end{proposition}
\begin{proof}
The projection respects source and target idempotents and extends
multiplicatively.
At a retained vertex, every term involving a discarded arrow vanishes.
The other terms are precisely $dt_i$ in the smaller algebra.
All arrow differentials vanish, so the projection is a chain map.
It preserves the path-length filtration and is continuous.
For the final assertion, start with one vertex $i$ and no arrows,
where $dt_i=0$.
Add a vertex $j$ and an arrow $a:i\to j$.
In the larger algebra $dt_i=aa^*$, a nonzero length-two path.
Thus the graded inclusion is not a dg map.
\end{proof}

The algebraic Ginzburg algebra for this finite quiver is homologically
smooth and three-Calabi--Yau \cite[Theorem~6.3]{KellerCY}. Its cyclic
class comes from the potential; zero potential satisfies the
negative-cyclic lift condition \cite{KellerCYErratum}.
The path-length-completed algebra is topologically homologically
smooth and three-Calabi--Yau in the pseudocompact bimodule category
\cite[Theorem~A.17]{KellerYang}. Its enveloping tensor product and
inverse-dualizing complex are continuous. This identifies the
finite completed model without replacing those tensors by ordinary
algebraic ones.
The displayed projections specify an inverse system for compatible
zero-potential windows.
A Calabi--Yau structure on a completed limit requires its own
nondegeneracy and continuity argument.
Comparisons with Hall categories and quantized Borcherds algebras must
also preserve orientations, root relations, and the paired Hopf data.

\section{The protected input}
The K3 elliptic genus is $2\phi_{0,1}$.
Use the normalized denominator $D_5:=\Dmon$ of Volume~III. The Borcherds half determines this weight-five form and its
Borcherds Lie superalgebra $\mathfrak g_\Delta$.
Fix its positive root spaces with their parity and proper height.
The positive algebra and its root Fock space are defined from this Lie
superalgebra, independently of a compact Hall realization.

For each finite zero-potential quiver window $Q_{\Delta,H}$, form the
Ginzburg algebra $\Gamma_{\Delta,H}$.
Restriction to full subquivers gives the projections of
Proposition~\ref{cyiv:ginzburg-projection}.
Write $\Gamma_\Delta$ for this specified inverse system.
A nondegenerate continuous Calabi--Yau structure on a completed limit
is additional data.
An ordinary Hall category $\mathcal B_\Delta$ must carry its chosen
quiver coefficients and extension correspondences.
Its identification with a root presentation requires the corresponding
Hall theorem on that precise datum.

\begin{definition}[Root-datum realizations]
A protected root-datum realization records
\[
 \mathfrak P(K3\times E)=
 (\Gamma_\Delta,\mathcal B_\Delta,\mathcal H_{\mathrm{1p}},
 \mathcal F^{\mathrm{hyb}}_{\Delta,E},\widehat U_q(\mathfrak g_\Delta),
 \mathbb A_Z).
\]
Here $\mathcal H_{\mathrm{1p}}=\bigoplus_{\alpha>0}(\mathfrak g_\Delta)_\alpha$
with its parity.
The current-factorization, paired quantum-group, and Fredholm entries
retain their own coefficient categories, completions, and defining
hypotheses.
Maps among these entries are specified realization maps, not consequences
of using the same root labels.
\end{definition}

The enveloping and Chevalley constructions are functorial under the
primitive maps of Theorem~\ref{cyiv:primitive-comparison}.
The root Fock trace is computed from the actual graded root spaces.
A Fredholm realization adds a family of operators on a specified domain
and the convergence and orientation data identifying its Weyl-shifted Pfaffian with
the denominator.
These are the analytic inputs to the square-root comparison below.

\begin{definition}[Integral protected realization]
Let $\mathfrak p_X^+$ be a BPS primitive Lie algebra supplied by a
selected reduced cohomological Hall realization, with an actual BPS
comparison or coproduct.  The integral protected realization is
\[
 \mathfrak P_{\mathrm{int}}(X,\Delta)
 =\left(
   \Hcal_X^{E,\mathrm{mot}},
   \mathfrak P(K3\times E),
   \mathfrak b_{X\boxtimes\Delta},
   r_X,r_\Delta,
   I_{X,\Pi}^E
 \right).
\]
The motivic component and its regular reduced character are defined as above; the root realizations retain their independent hypotheses.  The additional BPS or coproduct construction supplies
$\mathfrak p_X^+$; compatible additive Cartan weights then give the
amalgamated component.
\end{definition}

\begin{proposition}[Primitive graph comparison]
Suppose the two primitive Borels, their common Cartan action, and their
completed amalgam are supplied as above.
Their Cartan retractions induce retractions from the amalgam to each
factor.
For each specified primitive map $\varphi$ preserving Cartan action,
its graph quotient is canonically the geometric Borel.
\end{proposition}
\begin{proof}
The identity on one Borel and the Cartan projection of the other agree
on the common Cartan, so the amalgam gives the retraction.
The quotient statement is Proposition~\ref{cyiv:graph-quotient}.
\end{proof}

The motivic integration map and primitive comparison are different maps.
Strong orientation gives the coefficient coherence of a cohomological
Hall construction under its stated geometric hypotheses.
Extracting a primitive Lie algebra additionally requires the appropriate
BPS or coproduct construction.

\section{Cyclic open-string operations}
Let $A$ be a finite-dimensional, uncurved cyclic $A_\infty$ algebra
over $\mathbb C$, with a nondegenerate pairing of degree $-3$.
Use the cohomological suspension $s:A\to A[1]$ of degree $-1$.
The suspended operations $b_n:(A[1])^{\otimes n}\to A[1]$ have degree one.
Our sign convention is the coderivation convention $b^2=0$ on the
tensor coalgebra, with the Koszul rule for moving homogeneous inputs.
Write $\omega$ for the resulting degree-$-1$ symplectic form on $A[1]$.
Fix the Hamiltonian convention $\iota_{X_F}\omega=dF$ and
$[X_F,X_G]=X_{\{F,G\}}$.

For the universal degree-zero coordinate $x$ on $A[1]$, put
\[
 S(x)=\sum_{n\ge1}\frac1{n+1}\omega(x,b_n(x^{\otimes n})).
\]
Cyclic invariance is read with the suspended Koszul signs.
All functions are formal power series at the origin.
An equivalent unsuspended convention writes this as
$\frac12\langle a,m_1a\rangle+
\sum_{n\ge2}\frac1{n+1}\langle a,m_n(a^{\otimes n})\rangle$.

\begin{proposition}[Cyclic master equation]\label{cyiv:cyclic-master}
The action satisfies $\{S,S\}=0$.
\end{proposition}
\begin{proof}
Cyclicity makes the $n+1$ derivatives of the arity-$n+1$ term equal,
with the suspended Koszul signs.
Thus $X_S(x)=\sum_{n\ge1}b_n(x^{\otimes n})$.
The coderivation identity implies $X_S^2=0$.
The Hamiltonian identity gives $X_{\{S,S\}}=2X_S^2=0$.
Nondegeneracy of $\omega$ makes $\{S,S\}$ constant.
Uncurvedness removes $b_0$, so $S$ starts in order two and its bracket
has no constant term.
The constant is therefore zero.
\end{proof}

A finite-height Ginzburg realization requires such a cyclic model on
the selected proper objects before this scalar action is applied.
A Hall realization additionally requires a moduli problem with
critical coefficient complexes.
For an extension $A_0\to F\to B_0$, put
$L(A_0,B_0)=\det R\operatorname{Hom}(A_0,B_0)$.
Three-Calabi--Yau duality gives
\[
 R\operatorname{Hom}(B_0,A_0)
 \simeq R\operatorname{Hom}(A_0,B_0)^\vee[-3].
\]
Applying the determinant functor to the two derived Hom triangles gives
\[
 \det R\operatorname{Hom}(F,F)
 \cong \det R\operatorname{Hom}(A_0,A_0)
 \otimes\det R\operatorname{Hom}(B_0,B_0)
 \otimes L(A_0,B_0)^{\otimes2}.
\]
Indeed, the four subquotients are the two diagonal Hom complexes and
the two cross complexes.
Dualization and the odd shift cancel on their determinant lines,
so the two cross determinants coincide.
Orientation lines require square-root isomorphisms along this formula,
compatible on triple flags and with any quotient descent.
The cyclic master equation does not select these isomorphisms.

\section{An algebraic open--closed comparison}
A finite-dimensional ordinary Lie algebra $\mathfrak l$ gives a
precise coefficient model for an enveloping-algebra centre.
Let $U=U(\mathfrak l)$ over $\mathbb C$.
The adjoint module $U^{\mathrm{ad}}$ has action
$x\cdot u=xu-ux$.
Its Chevalley cochains are
\[
 C^n(\mathfrak l,U^{\mathrm{ad}})
 =\operatorname{Hom}_{\mathbb C}(\Lambda^n\mathfrak l,U),
\]
with differential
\[
\begin{split}
 (df)(x_1,\ldots,x_{n+1})
 &=\sum_i(-1)^{i+1}[x_i,f(x_1,\widehat x_i,\ldots,x_{n+1})]\\
 &\quad+\sum_{i<j}(-1)^{i+j}
 f([x_i,x_j],x_1,\widehat x_i,\widehat x_j,\ldots,x_{n+1}).
\end{split}
\]
The second line places the bracket first and retains the order of the
remaining arguments.

\begin{theorem}[Hochschild--Chevalley comparison]\label{cyiv:hh-ce}
Restriction followed by antisymmetrization defines a quasi-isomorphism
\[
 C^\bullet_{\mathrm{HH}}(U,U)
 \longrightarrow C^\bullet(\mathfrak l,U^{\mathrm{ad}}),
 \qquad
 \phi\longmapsto\left[(x_1,\ldots,x_n)\mapsto
 \sum_{\sigma\in S_n}\operatorname{sgn}(\sigma)
 \phi(x_{\sigma(1)},\ldots,x_{\sigma(n)})\right].
\]
This is a comparison of ordinary cochain complexes.
\end{theorem}
\begin{proof}
The Chevalley bimodule resolution has terms
$K_n=U\otimes\Lambda^n\mathfrak l\otimes U$.
Its differential has the left-minus-right multiplication terms and
the bracket terms with the signs dual to the displayed differential.
Filter $K_n$ by total PBW degree, counting each exterior generator
as degree one.
The multiplication terms preserve this degree and bracket terms lower it.
The associated graded is the Koszul resolution of the diagonal in
$\operatorname{Sym}(\mathfrak l)\otimes\operatorname{Sym}(\mathfrak l)$,
with regular sequence $x_i\otimes1-1\otimes x_i$.
That Koszul complex is exact away from degree zero.
The filtration is bounded below and exhaustive, and each element has
finite filtration degree, so lifting a leading term proves exactness
of $K_\bullet$ by induction.
Its augmentation has image $U$.
Thus $K_\bullet$ is a free bimodule resolution.

Antisymmetrization gives a map from $K_\bullet$ to the ordinary bar
resolution, lifting the identity of $U$.
Adjacent products in its bar differential pair to commutators;
$xy-yx=[x,y]$ gives exactly the bracket terms of $K_\bullet$.
The outer products give its left-minus-right terms.
A comparison between projective resolutions lifting the identity is
a homotopy equivalence.
Applying $\operatorname{Hom}_{U\otimes U^{\mathrm{op}}}(-,U)$ gives
the stated formula and quasi-isomorphism.
\end{proof}

This is the classical Cartan--Eilenberg comparison.
For a finite Lie superalgebra, the same construction uses the
super-exterior powers and permutation Koszul signs.
Odd generators produce symmetric exterior directions, so the resolution
can be unbounded; its algebraic direct sums and cochain products remain
part of the convention.
The PBW and super-Koszul proof applies with those conventions.

A physical mixed holomorphic--topological realization requires a map
from its independently defined bulk observables to this cochain object.
The theorem specifies the adjoint coefficients and the actual algebraic
comparison.
An $E_2$ refinement, a current-factorization realization, and compatibility
with root-height transitions require maps preserving their additional
operations.
A quotient of Lie algebras alone does not provide a covariant map of
Hochschild centres.

\section{Primitive comparisons and their algebraic extensions}
Let $\mathfrak n_\Delta$ be the positive Borcherds Lie superalgebra.
Choose homogeneous parity-preserving images of its simple generators
in a charge-graded Lie superalgebra $\mathfrak p$.
Assume these images satisfy every defining relation.

\begin{theorem}[Extension of a primitive comparison]\label{cyiv:primitive-comparison}
The chosen images determine a unique Lie-superalgebra map
$f:\mathfrak n_\Delta\to\mathfrak p$.
It induces
\[
 U(f)(x_1\cdots x_r)=f(x_1)\cdots f(x_r),
 \qquad C_*(f)=\operatorname{Sym}^c(f[1]).
\]
These are respectively an enveloping-algebra map and a Chevalley
chain-coalgebra map.
Filtration-preserving maps extend to the specified inverse-limit
completions.
If $f$ is surjective and both algebras have finite charge components
and proper positive height, equality of their parity-refined PBW
characters implies that $f$ and $U(f)$ are isomorphisms.
\end{theorem}
\begin{proof}
The defining presentation gives $f$.
Its tensor-algebra extension carries
$xy-(-1)^{|x||y|}yx-[x,y]$ to the corresponding relation.
It therefore descends to $U(f)$.
The shifted symmetric-coalgebra map commutes with the Chevalley
differential because $f$ preserves brackets.
Compatible maps on filtration quotients give maps on their inverse limits.
Surjectivity of $f$ gives surjectivity of $U(f)$.
Character equality makes each finite charge component an isomorphism,
separately in both parities.
The PBW embeddings then imply injectivity of $f$.
\end{proof}

At level zero the current Lie conformal algebra is
$\mathbb C[\partial]\otimes\mathfrak n_\Delta$ with
$[x_\lambda y]=[x,y]$ on constant generators.
The map $\partial^j x\mapsto\partial^j f(x)$ preserves this bracket
and sesquilinearity.
A nonzero central level additionally requires preservation of the
specified invariant pairing.
Hall convolution, Calabi--Yau structures, quantum coproducts, and
determinant lines require their own comparison data.

\section{The exact scalar and first-order observables}
Set
\[
 \DenIg(Z)=D_5(2Z),
 \qquad \jmath(Z)=Z/2,
 \qquad \jmath^*\DenIg=D_5.
\]
The one-particle root spaces give
\[
 \sch\Fock(\gDelta_+)
 =\prod_{\alpha>0}(1-e^{-\alpha})^{-\sdim\gDelta_\alpha}.
\]
Write $P(Z)=\prod_{\alpha>0}(1-e^{-\alpha(Z)})^{\sdim\gDelta_\alpha}$.
The normalized root operator gives, on its analytic domain,
\[
 \operatorname{sPf}_F(\mathbb A_Z)=P(Z),\qquad
 e^{-\rho}\operatorname{sPf}_F(\mathbb A_Z)=D_5(2Z).
\]
The primitive reduced curve-counting theorem gives
\[
 \mathcal Z^{\red}_{K3\times E}=-(\jmath^*\DenIg)^{-2}=-\chi_{10}^{-1}.
\]

\begin{theorem}[Protected square-root theorem]
Assume the stated root-space Fredholm comparison with its orientation,
Weyl shift, and convergence domain.
Within that realization, the denominator section $\DenIg$ is simultaneously:
\begin{enumerate}
\item the Weyl denominator of $\gDelta$;
\item the Weyl-shifted Euler character of the bar--Chevalley complex;
\item the Weyl-shifted oriented Fredholm Pfaffian of $\mathbb A_Z$;
\item the Weyl-shifted PBW determinant of the parity-refined root superspace $\mathcal H_{\mathrm{1p}}$;
\item the automorphic denominator whose half-coordinate pullback squares to the reduced scalar denominator.
\end{enumerate}
All five identifications preserve root degree. Across Weyl chambers
they preserve transport when the corresponding determinant-line
isomorphisms and orientation choices are included in the datum.
\end{theorem}
\begin{proof}
The denominator and bar--Chevalley theorems give $e^{-\rho}P$.
The Fredholm theorem gives $P$, and multiplication by the same Weyl
factor gives the third identification.  PBW for the Borcherds Lie superalgebra identifies the root superspace determinant, giving the fourth.  The half-coordinate identity $\jmath^*\DenIg=D_5$, the theta normalization, and the reduced curve-counting theorem give the fifth.
\end{proof}

\section{Physical meaning}
The root-datum realizations retain several distinct constructions.
The finite Ginzburg algebras and their projections form an inverse system.
A completed three-Calabi--Yau sector requires a continuous nondegenerate
Calabi--Yau structure and compatible cyclic models on the selected proper objects.
The quiver Hall category retains its chosen coefficients and extension
correspondences. A Hall theorem on that datum must identify its root presentation.
The Borcherds root superspace defines the root-graded Fock space and its
supercharacter. Its identification with physical one-particle states
requires a parity-preserving graded comparison intertwining the primitive actions.

A realization on the elliptic curve must construct the current object
with its local and wrapped supports, collision operations, and comparison maps.
The quantum Borcherds algebra, with its specified pairing and completion,
acts as a line-operator symmetry only through a representation preserving
fusion and exchange. The root operator gives the Fredholm product on its
specified convergence domain. A physical determinant interpretation requires
an operator comparison and an orientation-preserving identification of the
determinant lines and sections. The Weyl-normalized doubled Fock character gives the inverse scalar denominator.
An operator trace additionally requires its state space and trace domain.

Simple-root states satisfying the defining relations give a primitive Lie map.
The explicit enveloping and Chevalley extensions preserve products and brackets.
A compact string realization must additionally construct its brane functor,
Hall orientation maps, paired quantum coproduct, and trace-line comparison.

\section{Borcherds products as characters}
Let $\mathfrak n_+$ have finite-dimensional charge spaces and a
proper positive height. The alternating supercharacter of its Chevalley
chains, using homological sign and intrinsic parity, is
\[
 \chi_{\mathrm{s}}C_*(\mathfrak n_+)
 =\prod_{\alpha>0}(1-e^{-\alpha})^{\sdim(\mathfrak n_+)_\alpha}.
\]
A Borcherds product becomes the Weyl-shifted character of $C_*(\mathfrak n_+)$. This construction produces a character of an existing algebra. Hall geometry supplies a second algebra, and a primitive comparison theorem identifies the two.

\section{The universal Borcherds weight formula}
For a weight-zero weak Jacobi form of positive integral index satisfying the integrality hypotheses of the multiplicative Borcherds lift,
\[
 \varphi(\tau,z)=\sum c(n,\ell)q^nr^\ell,
\]
the Borcherds lift has weight
\[
 \wt(\operatorname{Borch}(\varphi))=\frac{c(0,0)}2.
\]
For the K3 half-index, $c(0,0)=10$ and the weight is $5$. This formula lives entirely on the automorphic layer. Hodge Euler characteristics, Heisenberg ranks, central charges, and BV anomalies are evaluated by their own traces.

\section{The Gritsenko--Cl\'ery atlas}
The eight diagonal-divisor forms are
\[
 \Delta_5,\Delta_2,\Delta_1,\Delta_{1/2},
 \nabla_3,\nabla_2,\nabla_{3/2},Q_1
\]
with weights
\[
 5,2,1,\frac12,3,2,\frac32,1.
\]
Each form is an automorphic product on its own paramodular or Hecke congruence group with its own character or multiplier \cite{GC}. A geometric realization is specified row by row by a charge lattice, Hall problem, trace, and primitive algebra comparison.

\section{Scalarization hierarchy}
Once the displayed comparison maps are constructed, one can pass
through the following recipients:
\[
\begin{aligned}
 \text{factorization or Hall object}
 &\longrightarrow \text{primitive Lie algebra}\\
 &\longrightarrow \text{PBW character}
 \longrightarrow \text{scalar partition function}.
\end{aligned}
\]
Primitive extraction requires a coproduct or a BPS comparison; the
last arrow requires a state, trace and normalization. These are not
automatic arrows for an arbitrary factorization algebra. The first object contains supports, extension maps, and brackets. The second retains the infinitesimal algebra. The third retains signed graded dimensions. The fourth adds a trace and normalization. Upward reconstruction is supplied by the corresponding geometric operations and comparison morphisms.

\chapter{Cotangent theories on K3 times E}

\section{Holomorphic cotangent BF theory on \texorpdfstring{$K3\times E$}{K3 x E}}
Let $S$ be a K3 surface with holomorphic symplectic form $\sigma_S$, let $E$ be an elliptic curve with holomorphic form $\dd z$, and put
\[
 X=S\times E,
 \qquad
 \Omega_X=\sigma_S\wedge\dd z.
\]
For a finite proper-height quotient $\nIg^{(H)}$ of the positive Igusa Borcherds algebra, take fields
\[
 A\in\Omega^{0,\bullet}(X,\nIg^{(H)})[1],
 \qquad
 B\in\Omega^{0,\bullet}(X,(\nIg^{(H)})^\vee)[1].
\]
The invariant evaluation pairing and $\Omega_X$ define the BV action
\begin{equation}\label{eq:vol4-compact-hbf}
 S_H(A,B)=
 \int_X\Omega_X\wedge
 \angles{B,\dbar A+\frac12[A,A]}.
\end{equation}
The two shifts are both $[1]$ in complex dimension three. After
integration against $\Omega_X$, evaluation has degree $-1$ on these
fields. This is the holomorphic cotangent theory of the dg Lie
superalgebra of $\nIg^{(H)}$-valued Dolbeault forms.

\begin{theorem}[Classical compact-background source]
The action \eqref{eq:vol4-compact-hbf} satisfies the classical master equation and defines a holomorphic factorization algebra on the compact Calabi--Yau threefold $X$.
\end{theorem}
\begin{proof}
The shifted cotangent pairing is invariant.  The classical master equation is the Jacobi identity of $\nIg^{(H)}$ together with $\dbar^2=0$ and the derivation property of $\dbar$.  Local observables of an elliptic cyclic $L_\infty$ algebra form a factorization algebra by the standard BV construction.
\end{proof}

\section{Loop weights and the quantum master equation}
Every interaction vertex in \eqref{eq:vol4-compact-hbf} has one $B$-leg and two $A$-legs.  The propagator pairs an $A$-leg with a $B$-leg.

\begin{lemma}[Loop bound]
A connected graph with at least one interaction vertex and the stated propagator has loop number at most one.
\end{lemma}
\begin{proof}
A graph with $V$ interaction vertices has exactly $V$ internal or external $B$-half-edges.  Every internal edge consumes one of them, hence $E\le V$.  The loop number is $E-V+1\le1$.
\end{proof}

The positive algebra is graded by strictly positive root height.  For homogeneous $x\in\nIg^{(H)}$, the operator $\ad_x$ raises height.  Its coadjoint operator lowers dual height.  A product of adjoint operators labelled by positive roots has total nonzero height shift and therefore has zero supertrace.

\begin{proposition}[Vanishing of one-loop Lie weights]\label{cyiv:loop-weights}
For every finite positive-height quotient, the Lie weight of each
one-loop graph with at least one interaction vertex vanishes.
\end{proposition}
\begin{proof}
A one-loop graph has $E=V$.
Every $B$-leg is consumed by an internal edge, so its external legs are
$A$-legs.
Its Lie weight is a cyclic supertrace of adjoint operators, or the
dual supertrace with the corresponding Koszul signs.
Each external label has positive root height.
The product therefore strictly raises height on the finite-dimensional
adjoint representation and has zero supertrace.
The same argument applies to the dual representation.
\end{proof}

This vanishing removes the one-loop Lie coefficients in a perturbative
construction with the stated propagator and vertex types.
A renormalized quantum factorization algebra still requires admissible
analytic kernels, local extensions, the quantum master equation, and
compatible height transitions on the chosen observable spaces.
These are the remaining quantum realization data. A quotient of
positive Lie algebras induces an inclusion on linear duals. Thus the
$A$- and $B$-field directions do not automatically give a same-direction
tower of cotangent theories or quantum observables.

\section{Classical reduction along K3}
Choose a K\"ahler metric on $S$ and normalize a closed harmonic
$(0,2)$-form $\eta$ by $\int_S\sigma_S\wedge\eta=1$.
Let $K=\mathbb C[\eta]/(\eta^2)$ with $|\eta|=2$ and zero differential.
There is a unital dg-algebra map
\[
 i:K\longrightarrow\Omega^{0,*}(S),\qquad \eta\longmapsto\eta.
\]
Its image is closed under multiplication because $(0,4)$-forms vanish.
The Dolbeault cohomology of a K3 surface is one-dimensional in degrees
zero and two and zero in degree one, so $i$ is a quasi-isomorphism.
It preserves the trace defined by integration against $\sigma_S$.

Let $\mathfrak l$ be a finite-dimensional Lie superalgebra, and write
$\mathfrak d=\mathfrak l\ltimes\mathfrak l^*$ for its cotangent Lie
superalgebra, with its invariant graded evaluation pairing.
For an open $U\subset E$, tensor with $\Omega^{0,*}(U)$, using the
completed smooth tensor product in the $S$ direction.
The field comparison is
\[
 K\otimes\Omega^{0,*}(U)\otimes\mathfrak d
 \xrightarrow{i\otimes\id}
 \Omega^{0,*}(S\times U)\otimes\mathfrak d.
\]
The bracket is multiplication of forms followed by the Lie bracket,
with Koszul signs.

\begin{proposition}[Classical field comparison]\label{cyiv:k3-classical-reduction}
The displayed map is a trace-preserving dg Lie quasi-isomorphism,
compatible with restriction in $U$.
Its induced minimal operations in the K3 factor have no higher products.
\end{proposition}
\begin{proof}
Multiplicativity of $i$ proves bracket compatibility.
Hodge projection $p$ and $h=\bar\partial^*G$ on the compact surface satisfy
$pi=\id$ and $\bar\partial h+h\bar\partial=\id-ip$.
Tensoring these continuous operators with forms on $U$ gives the same
contraction, with the tensor-product signs.
Thus $i\otimes\id$ is a quasi-isomorphism.
Restriction acts only on the $U$ factor and commutes with it.
The normalized trace on $K$ is the restriction of the Dolbeault trace.
Choose the contraction with $hi=0$.
In every higher transferred tree, the first product of inputs from
$i(K)$ remains in $i(K)$ and the next internal homotopy kills it.
Thus every higher transferred product vanishes.
\end{proof}

The coefficient algebra $K$ retains both the degree-zero and degree-two
K3 modes.
The proposition compares classical field complexes and their brackets.
Constructing an equivalence of quantum observable factorization algebras
requires compatible quantization, support conventions, and the induced
maps on observable completions.
The ordinary Hochschild comparison of Theorem~\ref{cyiv:hh-ce} is a
separate algebraic calculation, not that quantum field map.

\section{Protected states and the Igusa trace}
The formal root Fock object is
\[
 \widehat\Sym(\mathcal H_{\mathrm{1p},\bar0})
 \widehat\tensor\widehat\Lambda(\mathcal H_{\mathrm{1p},\bar1}),
 \qquad \mathcal H_{\mathrm{1p}}=\bigoplus_{\alpha>0}\gIg_\alpha,
\]
completed in proper positive height. Its character is a formal product.
For an operator trace, give each finite root space a Hilbert norm, take
their Hilbert direct sum, and form the bosonic and fermionic Hilbert
Fock spaces. Put $q_\alpha=e^{-\alpha(Z)}$ and let $e^{-H_Z}$ act
diagonally on one-particle states.

\begin{theorem}[Protected Hamiltonian trace]
Assume
\[
 \sum_{\alpha>0}\dim(\gIg_\alpha)|q_\alpha|<\infty,
 \qquad \sup_{\alpha>0}|q_\alpha|<1.
\]
Then its second quantization is trace class on the Hilbert Fock space,
and its supertrace is
\[
 \operatorname{str}(e^{-H_Z})
 =\prod_{\alpha>0}(1-q_\alpha)^{-\sdim\gIg_\alpha}=P(Z)^{-1}.
\]
The all-parity estimates of Volume~III give a nonempty domain satisfying
these hypotheses. With $\DenIg=e^{-\rho}P$, the primitive reduced
scalar is
\[
 -\bigl(\jmath^*\DenIg\bigr)^{-2}=-\chi_{10}^{-1},
 \qquad \jmath(Z)=Z/2,
\]
in the Oberdieck--Pixton normalization and primitive K3 curve sector.
\end{theorem}
\begin{proof}
The absolute trace of each even oscillator is
$(1-|q_\alpha|)^{-1}$, while that of each odd oscillator is
$1+|q_\alpha|$. The two hypotheses make the product of these absolute
traces finite, by the logarithm bounds on a disk of radius less than
one. Thus tensoring over the root basis gives a trace-class operator.
The supertrace replaces the odd factor by $1-q_\alpha$, giving
$P^{-1}$. The denominator identity is $\DenIg=e^{-\rho}P$.
Its half-coordinate pullback is $D_5$, whose square is $\chi_{10}$;
the primitive reduced curve-counting theorem supplies the final minus
sign \cite{OP}.
\end{proof}

This analytic trace is formed on the Hilbert completion just specified.
The formal product completion is a different state object. Identifying
either with a compact physical state space requires a representation
and trace comparison.

\section{The compact-source comparison}
The compact threefold provides the classical cotangent field theory
and the dg Lie comparison of Proposition~\ref{cyiv:k3-classical-reduction}.
The root datum separately gives the Borcherds algebra, its enveloping
algebra, and the root-graded Fock space.
The latter has the trace computed above.
Proposition~\ref{cyiv:loop-weights} supplies the finite one-loop Lie
cancellation, while Theorem~\ref{cyiv:hh-ce} computes the ordinary
algebraic centre model with adjoint coefficients.

\begin{problem}[Geometric compact comparison]\label{cyiv:compact-comparison}
Construct objects $E_i$ representing the simple branes and maps on their
derived morphism complexes that preserve units, composition, and higher
composition homotopies.
Realize these maps on extension stacks with the reduced critical
coefficients and orientation isomorphisms on triple flags.
Construct compatible tensor and unit constraints on the resulting
boundary functor.
Compare the chosen quantum bulk action with the independently defined
boundary centre, including support, collisions, and completion.
Finally construct a trace and a map of its value line to the automorphic
line in the same charge normalization.
\end{problem}

A primitive map satisfying the root relations gives the enveloping and
Chevalley maps of Theorem~\ref{cyiv:primitive-comparison}.
A strong monoidal equivalence transports categorical centres.
Neither construction determines the brane functor or its cyclic trace.
These are the additional maps required for compact geometric
reconstruction.

\section{Holomorphic--topological boundary theories}
A three-dimensional holomorphic--topological theory on $C\times\R$ has local observables that are holomorphic along $C$ and locally constant along $\R$. A factorizing boundary family with compatible composition gives the
ordered algebra constructed above. Local constancy and coherent
operations in two real directions give an $\Etwo$ structure when
topological factorization recognition applies. A meromorphic line
category also needs its holomorphic parameter and expansion data.
A coherent bulk action then factors through the stated centre.

The finite cotangent BF model gives the exact local calculation:
\[
 \Barc(U(\mathfrak l))\simeq C_*(\mathfrak l).
\]
For $\mathfrak l=\nDelta^{(H)}$, it computes the finite-height root
Chevalley character. Identifying it with boundary observables requires
the boundary construction; the automorphic denominator also carries
its Weyl factor.

\section{Six-dimensional hCS on \texorpdfstring{$K3\times E$}{K3 x E}}
Choose a holomorphic volume form $\Omega_S\wedge dz_E$. The hCS theory has classical fields
\[
 \Omega^{0,*}(S\times E,\mathfrak a)[1]
\]
and action $S_{\hCS}$. Its factorization observables are native to the threefold. Pushforward along $p_E$ gives a factorization algebra on $E$ whose coefficient complex contains the Dolbeault modes of $S$.

A protected reduction to a chiral algebra consists of a deformation retract of the K3 Dolbeault complex, compatibility of the transferred $L_\infty$ operations with the trace, a factorizing boundary or defect family with the coefficient comparisons above, and an admissible quantization solving the quantum master equation. These data produce the ordered chiral algebra to which the Hall/BPS comparison may be applied.

The Hodge equality $\chi(\Ocal_{S\times E})=0$ controls one compact supertrace. The gauge and gravitational anomalies are classes in the local BV deformation complex and are evaluated by their own characteristic forms.

\section{Omega backgrounds}
Let a torus act on the local geometry and suppose the cochain model
has operators $Q$ of degree $1$ and $\iota_V$ of degree $-1$, commuting with $\mathcal L_V$ and satisfying
\[
 Q^2=0,\qquad [Q,\iota_V]=\mathcal L_V,\qquad \iota_V^2=0.
\]
For an even parameter $u$ of cohomological degree two, put
\[
 Q_\Omega=Q+u\iota_V,\qquad Q_\Omega^2=u\mathcal L_V.
\]
It is a differential on the invariant subcomplex. A general
holomorphic differential does not carry these contraction operators
without additional data.

Equivariant localization in suitable toric or ADE models supplies
deformation parameters. A spectral $R$-matrix still requires its
correspondences and tensor identities. A compact K3 has no nonzero
holomorphic vector field: its symplectic form identifies such fields
with holomorphic one-forms, and $H^0(S,\Omega_S^1)=0$.
Consequently a localized model reaches a compact generic K3 only
through a specified deformation family and coefficient comparison,
with monodromy and wall-crossing compatibility.

\chapter{Anomalies and reconstruction}

\section{Anomalies as typed classes}
The constructions have the following obstruction recipients:
\begin{center}
\small
\begin{tabularx}{0.96\textwidth}{@{}>{\raggedright\arraybackslash}p{0.34\textwidth}>{\raggedright\arraybackslash}X@{}}
\toprule
Construction&Recipient\\\midrule
hCS quantization&local BV cochains\\
Critical CoHA orientation&Picard and gerbe theory of the critical stack\\
Hall coproduct&specified coefficient maps on flag correspondences\\
Projective conformal blocks&Atiyah algebra of the determinant line\\
Automorphic multiplier&character group of the modular cover\\
\bottomrule
\end{tabularx}
\end{center}
Some recipients are complexes, others are Picard groupoids, gerbes or
character groups. A comparison must first supply maps in the appropriate
categories; when an additive deformation-complex model exists, those
maps must identify the transported obstruction classes. This formulation turns numerical coincidences into geometric identities precisely when an intertwining map exists.

\section{The foundational theorem}
\begin{definition}[Complete Calabi--Yau quantum-group datum]
A complete comparison datum specifies independent categorical, field,
boundary, Hall, Hopf, and trace objects together with:
\begin{enumerate}
\item PTVV-representable moduli for the smooth proper Calabi--Yau category;
\item a renormalized BV quantization satisfying the quantum master equation and a factorizing family of local compact generators $b_I$;
\item a boundary reduction or proper locally trivial map to a curve satisfying factorization pushforward, together with the supplied curve-level faces and aligned distributive law;
\item a strong-oriented Hall coefficient theory whose pull--push maps satisfy base change and whose positive algebra is separated and complete;
\item paired negative and Cartan halves, a completed coproduct, a nondegenerate continuous Hopf pairing with its opposite-dual
convention, a continuous coevaluation or locally perfect dual pieces,
and a separated triangular completion with compatible Cartan reduction;
\item compatible open--closed and Hall--boundary morphisms;
\item a parity-refined primitive comparison for every automorphic identification.
\end{enumerate}
\end{definition}

\begin{theorem}[Calabi--Yau quantum-group assembly]
Fix a complete comparison datum as above. In particular it includes:
\begin{enumerate}
\item a smooth proper category with compatible cyclic Calabi--Yau data
and the representable moduli-of-objects hypotheses stated above;
\item a quantized cyclic elliptic BV theory and the compatible factorizing boundary or defect family;
\item a Hall extension correspondence with coefficient functor, base-change theorem, strong orientation data, and separated charge completion;
\item paired Hopf Borel halves with antipodes, compatible Cartan data,
a Hopf pairing, and the completed-double and coevaluation hypotheses;
\item every curve-level distributive, PBW, and analytic comparison used below.
\end{enumerate}
Then the datum determines functorially:
\begin{enumerate}
\item framed $\Etwo$ Hochschild cochains and the shifted-symplectic moduli structure;
\item the three curve-level algebraic faces and their three bars;
\item the positive Hall algebra and, when a BPS comparison or coproduct
is supplied, its specified primitive Lie algebra;
\item the completed quantum double and braided representation category;
\item every supplied comparison among categorical, field-theoretic, Hall, current, quantum, and automorphic realizations;
\item the Weyl-shifted PBW/Chevalley character of any parity-refined primitive Lie-superalgebra identification.
\end{enumerate}
\end{theorem}
\begin{proof}
Cyclic Deligne and shifted symplectic geometry give the first output \cite{BravRozenblyum,PTVV,BravDyckerhoff}.  Boundary factorization gives the curve faces, and Volume~I gives their bars and centres.  The Hall correspondence and its base-change theorem give the associative product; strong orientation supplies the coefficient coherence required in the $3$-Calabi--Yau setting \cite{KPS}.  The completed-double hypotheses give the double and its canonical
$R$-matrix on the stated module domain.  The comparison maps commute by hypothesis.  PBW and Chevalley antisymmetrization give the final character.  Each construction is natural under a morphism preserving the displayed data.
\end{proof}

The local quiver, equivariant surface, K3 correspondence, and Igusa Borcherds constructions have the carriers stated above.
A geometric $K3\times E$ comparison requires the primitive map and determinant-line map of Problem~\ref{cyiv:compact-comparison}.

\section{Landscape of established realizations}
\begin{center}
\footnotesize
\begin{tabularx}{0.99\textwidth}{@{}>{\raggedright\arraybackslash}p{0.15\textwidth}>{\raggedright\arraybackslash}p{0.25\textwidth}>{\raggedright\arraybackslash}X>{\raggedright\arraybackslash}p{0.20\textwidth}@{}}
\toprule
\textbf{Geometry}&\textbf{Positive algebra}&\textbf{Full or centred object}&\textbf{Character}\\
\midrule
$A_2$ quiver&$U_v^+(\mathfrak{sl}_3)$&$U_v(\mathfrak{sl}_3)$&PBW root product\\
$\C^3$&$Y^+(\widehat{\mathfrak{gl}}_1)$&$Y(\widehat{\mathfrak{gl}}_1)$&plane-partition and Fock traces\\
$\widetilde{\C^2/\Gamma}$&$Y^+_\infty(\widehat{\mathfrak g}_\Gamma)$&paired completion when the negative half and coproduct are supplied&surface Fock traces\\
$K3$&BPS envelope; $W^+(S)$ for zero-dimensional sheaves&global
$W^{(r)}(S)$ actions; surface bialgebra in its pro-coefficients&K3 Hall characters\\
$K3\times E$&Borcherds root envelope; geometric reduced character&$U_q(\gDelta)$ and the root-current model and classical BF theory&$-\ChiTen^{-1}$\\
\bottomrule
\end{tabularx}
\end{center}
The $A_2$ and $\C^3$ rows provide complete standard quantum-group models. The ADE row provides an established completed positive half together with its geometric current actions; a paired full double additionally requires all the Hopf, coevaluation and topology hypotheses stated above. The K3 row includes the global BPS algebra, the zero-dimensional
$W$-algebra and its associative double actions, and the pro-cohomological
surface coproduct with its functorial coefficient hypotheses. The final row records independently constructed root, character, and classical field objects.
Their geometric Hall realization requires the comparison in Problem~\ref{cyiv:compact-comparison}.

\section{Constructed comparisons and their geometric extensions}
The Hall flag argument constructs associativity once its coefficient
and orientation maps are defined.
The tensor-action formula constructs representation fusion from a
coassociative counital coproduct.
The matrix-coefficient formulas construct the coend of a supplied
fibre functor.
The primitive presentation gives enveloping and Chevalley maps.
The classical cyclic and K3 field calculations retain their actual
differentials, traces, and brackets.

The quantum field realization requires admissible quantization and
observable completions.
The compact Hall realization requires the brane and orientation maps
of Problem~\ref{cyiv:compact-comparison}.
The Fredholm and automorphic comparisons require their analytic domains
and value-line isomorphisms.
These constructions preserve the full geometric reconstruction problem
beyond its finite algebraic and classical calculations.

\section{Actions and comparison maps}
For a specified boundary algebra $A_b$, a coherent bulk action is a map
\[
 \beta:\Obs_{\mathrm{bulk}}|_\partial\longrightarrow Z_{\ord}(A_b).
\]
The target is the derived endomorphism object of the diagonal bimodule
in the chosen factorization coefficient category.
A Hall construction instead starts with
\[
 \M_\alpha\times\M_\beta\xleftarrow{p}\mathcal E_{\alpha,\beta}
 \xrightarrow{q}\M_{\alpha+\beta}.
\]
Critical coefficients, orientations, and admissible pull--push maps
give its product.
Framed moduli supply a representation after their action
correspondences and triple-flag identities have been constructed.
A paired Hopf structure gives $D(H^+,H^0,H^-)$.
Its coproduct defines tensor products of continuous modules, and a
convergent universal $R$-matrix gives exchange.

A strong monoidal equivalence between this module category and a
boundary category transports categorical centres.
It does not identify that categorical centre with the object
$Z_{\ord}(A_b)$ without a realization theorem.
A trace on a specified representation then gives an amplitude.
Its comparison with an automorphic section requires a map of value
lines with the same normalization.

\section{Comparison data for Calabi--Yau constructions}
Comparisons among Calabi--Yau categories, holomorphic field theory, chiral algebras, Hall algebras, quantum groups, and automorphic products require the following structures:
\begin{enumerate}
\item for a smooth proper $d$-Calabi--Yau category with compatible
nondegenerate left and right cyclic structures, the categorical outputs are framed $\Etwo$ Hochschild cochains and, when the moduli stack of objects is representable in the PTVV setting, $(2-d)$-shifted symplectic moduli; the shift $d$ does not change the Hochschild operadic level to $\En$;
\item a geometric BV theory gives factorization observables, and a
factorizing boundary family with coherent composition gives a
transverse ordered chiral algebra; associative chiral convolution is recorded as a separate supplied structure;
\item pushforward to a curve uses an actual map $p:X\to C$; formal cycle-specialization notation alone specifies no pushforward;
\item extension stacks, coefficient sheaves, orientations, and completions define the Hall algebra;
\item Hopf Borel halves, opposite dual convention, Cartan reduction,
pairing, coevaluation and topology define the stated double;
\item K3 Heisenberg, LLV, local ADE Yangians, the global surface
$W$-algebra and affinized BPS algebra have their stated carriers and
coefficient systems; constructed homomorphisms join their actions;
\item the graph of the $K3\times E$ Gram map becomes multiplicative in the polarized central extension;
\item point-local and wrapped sectors form a hybrid factorization problem;
\item the Igusa form is the Weyl-shifted PBW/Chevalley character of $\gDelta$, and a compact Hall identification requires the primitive and trace-line comparisons;
\item anomaly identities are transported through explicit maps between
their complexes, Picard objects or character groups.
\end{enumerate}
The categorical, geometric, Hall, and automorphic operations thus have separate carriers and comparison maps.

\section{Typed data table}
\begin{center}
\small
\begin{tabularx}{0.98\textwidth}{@{}>{\raggedright\arraybackslash}p{0.20\textwidth}>{\raggedright\arraybackslash}X>{\raggedright\arraybackslash}p{0.28\textwidth}@{}}
\toprule
\textbf{Object}&\textbf{Minimal data}&\textbf{Primary operation}\\
\midrule
$\CY$ category&compatible smooth/proper cyclic data&Hochschild and shifted-symplectic structures\\
Factorization algebra&local observables and Weiss descent&disjoint-open product\\
Transversely ordered chiral algebra&factorization coefficient with an internal $\Eone$ structure for $\tensor^!$&ordered fusion\\
Associative chiral algebra&an $\Eone$ structure for chiral convolution $\starCh$&nonlocal chiral composition\\
$\CoHA$&extension correspondence and coefficient theory&pull--push product\\
Bialgebra&compatible product and coassociative counital coproduct&fusion of representations\\
Double&Hopf halves, duality, Cartan reduction and convergent coevaluation&universal $R$-matrix\\
Borcherds product&Jacobi or vector-valued modular input&graded character\\
\bottomrule
\end{tabularx}
\end{center}

\section{Low-degree formulas}
Several small formulas distinguish the operations in this construction.
The K3 half-index begins
\[
 \phi_{0,1}(\tau,z)=r^{-1}+10+r+O(q),
\]
and its higher Fourier rows are obtained by the exact theta division
through $q^3$ in Volume~III. The constant coefficient $10$ gives weight $5$.
The Gram polarization satisfies
\[
 B_\Pi(a,b)+B_\Pi(a+b,d)
 =B_\Pi(b,d)+B_\Pi(a,b+d),
\]
which proves the extension law used above. The rational exchange matrix
$1-\hbar P/u$ satisfies the three-factor Yang--Baxter equation; its
normalization for two-sided unitarity is separate, as in Volume~I.
The two $A_2$ flag rows give the quantum Serre identity. Finally the
type-II Gram matrix in the basis $r_1,r_2,r_3$ and its simple reflections are
\[
 G=\begin{pmatrix}2&-2&-2\\-2&2&-2\\-2&-2&2\end{pmatrix},
 \qquad s_i(v)=v-(v,r_i)r_i.
\]
Since $(r_i,r_i)=2$, direct expansion gives $s_i^2=1$ and
$(s_i v,s_i w)=(v,w)$. Thus these reflections preserve the lattice
form; they do not identify it with the additive Cartan weights.

